\input{preamble}

% OK, start here.
%
\begin{document}

\title{Corps}


\maketitle

\phantomsection
\label{section-phantom}

\tableofcontents


\section{Introduction}
\label{section-introduction}

\noindent
Dans ce chapitre, nous étudierons la théorie des corps. Rappelons qu'un
{\it corps} est un anneau non nul dans lequel tout élément non nul est inversible.
De manière équivalente, les deux seuls idéaux d'un corps sont $(0)$ et $(1)$,
puisque tout élément non nul est une unité. Les corps constitueront donc les
cas les plus simples d'une grande partie de la théorie développée ultérieurement.

\medskip\noindent
La théorie des extensions de corps présente un aspect différent de l'algèbre
commutative classique puisque, par exemple, tout morphisme de corps est injectif.
Il s'avère néanmoins que des questions portant sur les anneaux peuvent souvent
se ramener à des questions sur les corps. Par exemple, tout anneau intègre se plonge dans un corps
(son corps des fractions), et à tout {\it anneau local} (c'est-à-dire un anneau possédant un unique
idéal maximal ; nous n'avons pas encore défini ce terme) est associé son
corps résiduel (c'est-à-dire son quotient par l'idéal maximal).
La connaissance des extensions de corps sera donc utile.




\section{Définitions de base}
\label{section-definitions}

\noindent
Comme nous avons placé ce chapitre avant celui consacré à l'algèbre
commutative, nous devons introduire ici certaines définitions élémentaires
avant de les étudier plus en détail dans les chapitres d'algèbre.

\begin{definition}
\label{definition-field}
Un {\it corps} est un anneau non nul dont tout élément non nul est inversible.
Étant donné un corps, un {\it sous-corps} est un sous-anneau qui est lui-même un corps.
\end{definition}

\noindent
Pour un corps $k$, nous notons $k^*$ le sous-ensemble $k \setminus \{0\}$.
Cette notation généralise la notation usuelle $R^*$ qui désigne le groupe des
éléments inversibles d'un anneau $R$.

\begin{definition}
\label{definition-domain}
Un {\it anneau intègre}, ou un {\it domaine d'intégrité}, est un anneau non nul dans lequel $0$
est le seul diviseur de zéro.
\end{definition}



\section{Exemples de corps}
\label{section-examples}

\noindent
Pour commencer, donnons plusieurs exemples de corps. Le lecteur se rappellera
que, si $R$ est un anneau et $I \subset R$ un idéal, alors $R/I$ est un corps
si et seulement si $I$ est un idéal maximal.

\begin{example}[Nombres rationnels]
\label{example-rational-numbers}
Les nombres rationnels forment un corps. On l'appelle le
{\it corps des nombres rationnels} et on le note $\mathbf{Q}$.
\end{example}

\begin{example}[Corps premiers]
\label{example-prime-field}
Si $p$ est un nombre premier, alors $\mathbf{Z}/(p)$ est un corps, noté
$\mathbf{F}_p$. En effet, $(p)$ est un
idéal maximal de $\mathbf{Z}$. Ainsi, un corps peut être fini : $\mathbf{F}_p$
contient $p$ éléments.
\end{example}

\begin{example}
\label{example-quotient-polymial-ring}
Dans un anneau principal, tout idéal engendré par un élément irréductible
est maximal. Or, si $k$ est un corps, l'anneau de polynômes $k[x]$ est un
anneau principal. Il s'ensuit que, si $P \in k[x]$ est un polynôme irréductible
(c'est-à-dire un polynôme non constant
qui n'admet pas de factorisation en facteurs de degrés plus petits), alors
$k[x]/(P)$ est un corps. Il contient naturellement une copie de $k$.
C'est une méthode très générale pour construire des corps. Par exemple, les
nombres complexes $\mathbf{C}$
peuvent être construits sous la forme $\mathbf{R}[x]/(x^2 + 1)$.
\end{example}

\begin{example}[Corps des fractions]
\label{example-quotient-field}
Rappelons qu'étant donné un anneau intègre $A$, il existe un plongement $A \to F$ dans un
corps $F$ construit à partir de $A$ exactement de la même manière que
$\mathbf{Q}$ est construit à partir de $\mathbf{Z}$. Formellement, les éléments
de $F$ sont des fractions $a/b$ (ou, plus précisément, leurs classes d'équivalence),
$a, b \in A$, $b \not = 0$. Comme d'habitude, $a/b = a'/b'$ si et seulement si $ab' = ba'$.
Le corps $F$ est appelé le {\it corps des fractions} de $A$ ; les calques
{\it corps quotient} et {\it corps de fractions} ne seront pas employés.
Le corps des fractions possède la propriété universelle suivante : pour tout
morphisme injectif d'anneaux $\varphi : A \to K$ vers un corps $K$, il existe un unique
morphisme $\psi : F \to K$ qui rende le diagramme
$$
\xymatrix{
F \ar[r]_\psi & K \\
A \ar[u] \ar[ru]_\varphi
}
$$
commutatif. En effet, la façon de définir un tel morphisme est claire : on pose
$\psi(a/b) = \varphi(a)\varphi(b)^{-1}$, l'injectivité de $\varphi$
assurant que $\varphi(b) \not = 0$ si $ b \not = 0$.
\end{example}

\begin{example}[Corps des fonctions rationnelles]
\label{example-field-of-rational-functions}
Si $k$ est un corps, on peut considérer le corps $k(x)$ des fonctions
rationnelles sur $k$. C'est le corps des fractions de l'anneau de polynômes
$k[x]$. Autrement dit, c'est l'ensemble des quotients $F/G$ avec
$F, G \in k[x]$, $G \not = 0$, muni de la relation d'équivalence évidente.
\end{example}

\begin{example}
\label{example-field-of-meromorphic-functions}
Soit $X$ une surface de Riemann. Notons $\mathbf{C}(X)$
l'ensemble des fonctions méromorphes sur $X$. Alors $\mathbf{C}(X)$ est un anneau pour
la multiplication et l'addition des fonctions. Il s'avère qu'en fait
$\mathbf{C}(X)$ est un corps. En effet, si une fonction non nulle $f(z)$ est
méromorphe, il en va de même de $1/f(z)$. Par exemple, soit $S^2$ la sphère
de Riemann ; l'analyse complexe nous apprend que l'anneau des fonctions méromorphes
$\mathbf{C}(S^2)$ est le corps des fonctions rationnelles $\mathbf{C}(z)$.
\end{example}



\section{Espaces vectoriels}
\label{section-vector-spaces}

\noindent
L'une des raisons pour lesquelles les corps sont si commodes est que la théorie des modules sur un corps
(c'est-à-dire des espaces vectoriels) est très simple.

\begin{lemma}
\label{lemma-vector-space-is-free}
Si $k$ est un corps, tout $k$-module est libre.
\end{lemma}

\begin{proof}
En effet, l'algèbre linéaire nous apprend qu'un $k$-module (c'est-à-dire un espace vectoriel)
$V$ possède une {\it base} $\mathcal{B} \subset V$, laquelle définit un isomorphisme
de l'espace vectoriel libre sur $\mathcal{B}$ vers $V$.
\end{proof}

\begin{lemma}
\label{lemma-field-semi-simple}
Toute suite exacte de modules sur un corps est scindée.
\end{lemma}

\begin{proof}
Cela résulte du Lemme \ref{lemma-vector-space-is-free}, puisque tout espace
vectoriel est un module projectif.
\end{proof}

\noindent
C'est une autre raison pour laquelle une grande partie de la théorie des chapitres ultérieurs dira
peu de choses sur les corps, puisque les modules s'y comportent de manière si simple.
Notons que le Lemme \ref{lemma-field-semi-simple} est un énoncé portant sur la
{\it catégorie} des $k$-modules (où $k$ est un corps), car la notion
d'exactitude est intrinsèquement de nature morphique, c'est-à-dire qu'elle ne fait intervenir que des notions
catégoriques, et peut en fait se formuler dans ce que l'on appelle une {\it catégorie abélienne}.

\medskip\noindent
Désormais, puisque l'étude des modules sur un corps relève de l'algèbre linéaire et
que la théorie des idéaux d'un corps présente peu d'intérêt, nous étudierons le véritable
objet de ce chapitre : les {\it extensions} de corps.


\section{La caractéristique d’un corps}
\label{section-more-fields}

\noindent
Dans la catégorie des anneaux, il existe un {\it objet initial} $\mathbf{Z}$ : tout
anneau $R$ reçoit un morphisme de $\mathbf{Z}$ d'une manière et d'une seule. Pour les corps,
il n'existe pas d'objet initial.
Il existe néanmoins une famille d'objets telle que tout corps reçoive d'une manière
unique un morphisme provenant d'exactement l'un d'entre eux, et aucun morphisme provenant des autres.

\medskip\noindent
Soit $F$ un corps. Considérons $F$ comme un anneau afin d'obtenir un morphisme d'anneaux
$f : \mathbf{Z} \to F$. L'image de ce morphisme d'anneaux est un anneau intègre
(en tant que sous-anneau d'un corps) ; le noyau de $f$ est donc un idéal premier
de $\mathbf{Z}$. Par conséquent, le noyau de $f$ est soit $(0)$, soit $(p)$ pour
un certain nombre premier $p$.

\medskip\noindent
Dans le premier cas, $f$ est injectif, et nous considérons alors
$\mathbf{Z}$ comme un sous-anneau de $F$. De plus, puisque tout
élément non nul de $F$ est inversible, il est légitime de
parler de $p/q \in F$ pour $p, q \in \mathbf{Z}$ avec $q \not = 0$.
Dans ce cas, nous pouvons donc considérer, et nous considérons, $\mathbf{Q}$ comme un sous-anneau de $F$.
On vérifie aisément que c'est alors le plus petit sous-corps de $F$.

\medskip\noindent
Dans le second cas, c'est-à-dire lorsque $\Ker(f) = (p)$, on voit que
$\mathbf{Z}/(p) = \mathbf{F}_p$ est un sous-anneau de $F$. C'est manifestement le
plus petit sous-corps de $F$.

\medskip\noindent
Ce raisonnement montre que tout corps contient un plus petit sous-corps,
qui est soit $\mathbf{Q}$, soit le corps fini $\mathbf{F}_p$ pour un certain
nombre premier $p$.

\begin{definition}
\label{definition-characteristic}
La {\it caractéristique} d'un corps $F$ est $0$ si
$\mathbf{Z} \subset F$, et est un nombre premier $p$ si $p = 0$ dans $F$.
Le {\it sous-corps premier de $F$} est le plus petit sous-corps de $F$ ;
il s'agit soit de $\mathbf{Q} \subset F$ si la caractéristique est nulle, soit de
$\mathbf{F}_p \subset F$ si la caractéristique est $p > 0$.
\end{definition}

\noindent
Il est facile de voir que, si $E \subset F$ est un sous-corps, la
caractéristique de $E$ est la même que celle de $F$.

\begin{example}
\label{example-characteristic}
La caractéristique de $\mathbf{F}_p$ est $p$, et celle de $\mathbf{Q}$ est $0$.
\end{example}


\section{Extensions de corps}
\label{section-extensions}

\noindent
En général, cependant, ce ne sont pas tant les corps eux-mêmes qui nous intéressent que leurs
{\it extensions}. Cette démarche est peut-être analogue à celle qui consiste à étudier non les anneaux,
mais les {\it algèbres} sur un anneau fixé.
Dans le cas des corps, la notion de ``corps sur un autre corps''
redonne simplement celle d'extension de corps, comme le montre le résultat suivant.

\begin{lemma}
\label{lemma-field-maps-injective}
Si $F$ est un corps et $R$ un anneau non nul, alors tout homomorphisme d'anneaux
$\varphi : F \to R$ est injectif.
\end{lemma}

\begin{proof}
En effet, soit $a \in \Ker(\varphi)$ un élément non nul. On a alors
$\varphi(1) = \varphi(a^{-1} a) = \varphi(a^{-1}) \varphi(a) = 0$.
Ainsi $1 = \varphi(1) = 0$, et $R$ est l'anneau nul.
\end{proof}

\begin{definition}
\label{definition-extension}
Si $F$ est un corps contenu dans un corps $E$, on dit que $E$ est
une {\it extension de corps} de $F$. Nous écrirons $E/F$ pour indiquer
que $E$ est une extension de $F$.
\end{definition}

\noindent
Ainsi, si $F, F'$ sont des corps et $F \to F'$ un homomorphisme d'anneaux quelconque, le
Lemme \ref{lemma-field-maps-injective} montre qu'il est injectif, et $F'$ peut être
considéré comme une extension de $F$, par un léger abus de langage. De façon équivalente,
une extension de corps de $F$ n'est autre qu'une $F$-algèbre qui est un corps.
La situation est entièrement différente pour les anneaux généraux, puisqu'un
homomorphisme d'anneaux n'est pas nécessairement injectif.

\medskip\noindent
Soit $k$ un corps. Il existe une {\it catégorie} des extensions de corps de $k$.
Un objet de cette catégorie est une extension $E/k$, c'est-à-dire un
morphisme de corps (nécessairement injectif)
$$
k \to E,
$$
tandis qu'un morphisme entre les extensions $E/k$ et $E'/k$ est un morphisme de
$k$-algèbres $E \to E'$ ; de façon équivalente, c'est un diagramme commutatif
$$
\xymatrix{
E \ar[rr] & & E' \\
& k \ar[ru] \ar[lu] &
}
$$
L'ensemble des morphismes $E \to E'$ dans la catégorie des extensions de $k$
sera noté $\Mor_k(E, E')$.

\begin{definition}
\label{definition-tower}
Une {\it tour} de corps $E_n/E_{n - 1}/\ldots/E_0$ consiste en une suite
d'extensions de corps
$E_n/E_{n - 1}$, $E_{n - 1}/E_{n - 2}$, $\ldots$, $E_1/E_0$.
\end{definition}

\noindent
Donnons quelques exemples d'extensions de corps.

\begin{example}
\label{example-monogenic-extension}
Soient $k$ un corps et $P \in k[x]$ un polynôme irréductible. Nous avons
vu que $k[x]/(P)$ est un corps (Exemple \ref{example-quotient-polymial-ring}).
Comme c'est aussi une $k$-algèbre de façon évidente, il s'agit d'une extension de $k$.
\end{example}

\begin{example}
\label{example-field-of-meromorphic-functions-extension-C}
Si $X$ est une surface de Riemann, le corps des fonctions méromorphes
$\mathbf{C}(X)$ (Exemple \ref{example-field-of-meromorphic-functions})
est une extension de corps de $\mathbf{C}$, car tout élément de $\mathbf{C}$
définit sur $X$ une fonction constante méromorphe --- et même holomorphe.
\end{example}

\noindent
Soit $F/k$ une extension de corps. Soit $S \subset F$ un sous-ensemble quelconque.
Il existe alors une {\it plus petite} sous-extension de $F$ (c'est-à-dire un sous-corps de
$F$ contenant $k$) qui contienne $S$. Pour le voir, considérons la famille des
sous-corps de $F $ contenant $S$ et $k$, puis prenons leur intersection ; on
vérifie qu'il s'agit d'un corps. Un argument classique montre en fait que
c'est l'ensemble des éléments de $F$ que l'on peut obtenir au moyen d'un nombre fini
d'opérations algébriques élémentaires (addition, multiplication, soustraction
et division) portant sur des éléments de $k$ et de $S$.

\begin{definition}
\label{definition-generated-by}
Soit $k$ un corps. Si $F/k$ est une extension de corps et
$S \subset F$, nous notons $k(S)$ le plus petit sous-corps de $F$
contenant $k$ et $S$. Nous dirons que $S$ {\it engendre
l'extension de corps} $k(S)/k$. Si $S = \{\alpha\}$ est un singleton, nous
écrirons $k(\alpha)$ au lieu de $k(\{\alpha\})$. On dit que $F/k$ est une
{\it extension de corps de type fini} s'il existe un
sous-ensemble fini $S \subset F$ tel que $F = k(S)$.
\end{definition}

\noindent
Par exemple, $\mathbf{C}$ est engendré par $i$ sur $\mathbf{R}$.

\begin{exercise}
\label{exercise-C-not-countably-generated}
Montrer que $\mathbf{C}$ ne possède pas de famille dénombrable de générateurs sur
$\mathbf{Q}$.
\end{exercise}

\noindent
Classifions maintenant les extensions engendrées par un élément.

\begin{lemma}[Classification des extensions simples]
\label{lemma-field-extension-generated-by-one-element}
Si une extension de corps $F/k$ est engendrée par un élément, alors elle est
$k$-isomorphe soit au corps des fonctions rationnelles $k(t)/k$, soit à l'une
des extensions $k[t]/(P)$, où $P \in k[t]$ est irréductible.
\end{lemma}

\noindent
Nous verrons que bon nombre des cas les plus importants d'extensions de corps sont
engendrés par un élément ; ce résultat est donc effectivement utile.

\begin{proof}
Soit $\alpha \in F$ tel que $F = k(\alpha)$ ; par hypothèse, un tel
$\alpha$ existe. Il existe un morphisme d'anneaux
$$
k[t] \to F
$$
qui envoie l'indéterminée $t$ sur $\alpha$. Son image est un anneau intègre, donc son
noyau est un idéal premier. Celui-ci est donc soit $(0)$, soit $(P)$ pour un $P \in k[t]$
irréductible.

\medskip\noindent
Si le noyau est $(P)$ pour un $P \in k[t]$ irréductible, le morphisme se factorise
par $k[t]/(P)$ et induit un morphisme de corps $k[t]/(P) \to F$. Puisque
l'image contient $\alpha$, on voit aisément que ce morphisme est surjectif, donc
est un isomorphisme. Dans ce cas, $k[t]/(P) \simeq F$.

\medskip\noindent
Si le noyau est trivial, on obtient une injection $k[t] \to F$.
On peut donc définir un morphisme du corps des fractions $k(t)$ dans $F$ ; à un
quotient $R(t)/Q(t)$ avec $R(t), Q(t) \in k[t]$, on associe
$R(\alpha)/Q(\alpha)$. L'hypothèse selon laquelle $k[t] \to F$ est injectif implique
que $Q(\alpha) \neq 0$, sauf si $Q$ est le polynôme nul.
Le corps des fractions de $k[t]$ est le corps des fonctions rationnelles $k(t)$ ; on obtient donc
un morphisme $k(t) \to F$
dont l'image contient $\alpha$. Il est par conséquent surjectif, donc est un isomorphisme.
\end{proof}

\begin{lemma}
\label{lemma-common-extension-field}
Soient $k$ un corps et $E/k$, $F/k$ des extensions de corps.
Il existe alors une extension de corps commune $M/k$, c'est-à-dire une extension
de corps telle qu'il existe des morphismes $E \to M$ et $F \to M$ d'extensions de $k$.
\end{lemma}

\begin{proof}
Nous ne démontrons ce résultat que lorsque $E$ est une extension de corps de type fini
de $k$ ; le cas général s'en déduit par un argument de type lemme de Zorn
(les détails sont omis).

\medskip\noindent
Supposons d'abord que $E$ soit une extension simple de $k$. D'après le
Lemme \ref{lemma-field-extension-generated-by-one-element},
cela signifie soit que $E = k(t)$ est le corps des fonctions rationnelles,
soit que $E = k[t]/(P)$ pour un certain polynôme irréductible $P \in k[t]$.
Dans le premier cas, prenons pour $M = F(t)$ le corps des fonctions rationnelles,
muni des morphismes évidents $E \to M$ et $F \to M$. Dans le second cas, nous
choisissons un facteur irréductible $Q$ de l'image de $P$ dans $F[t]$
et prenons $M = F[t]/(Q)$, muni des morphismes évidents $E \to M$ et $F \to M$.

\medskip\noindent
Si $E = k(\alpha_1, \ldots, \alpha_n)$, une récurrence sur $n$
permet de trouver une extension $M/k$ et des morphismes $F \to M$ et
$k(\alpha_1, \ldots, \alpha_{n - 1}) \to M$. En appliquant le cas simple
étudié au paragraphe précédent,
on trouve une extension $M'/k(\alpha_1, \ldots, \alpha_{n - 1})$
et des morphismes $M \to M'$ et $k(\alpha_1, \ldots, \alpha_n) \to M'$.
Alors $M'$, considéré comme une extension de $k$, convient.
\end{proof}




\section{Extensions finies}
\label{section-finite-extensions}

\noindent
Si $F/E$ est une extension de corps, alors $F$ est aussi, de manière évidente, un espace vectoriel
sur $E$ (l'action scalaire n'est autre que la multiplication dans $F$).

\begin{definition}
\label{definition-degree}
Soit $F/E$ une extension de corps. La dimension de $F$, considéré comme un
$E$-espace vectoriel, est appelée le {\it degré} de l'extension et est
notée $[F : E]$. Si $[F : E] < \infty$, on dit que $F$ est une extension
{\it finie} de $E$.
\end{definition}

\begin{example}
\label{example-C-over-R}
Le corps $\mathbf{C}$ est un espace vectoriel de dimension deux sur $\mathbf{R}$,
de base $1, i$. Ainsi $\mathbf{C}$ est une extension finie de $\mathbf{R}$
de degré 2.
\end{example}

\begin{lemma}
\label{lemma-finite-goes-up}
Soit $K/E/F$ une tour d'extensions algébriques de corps.
Si $K$ est fini sur $F$, alors $K$ est fini sur $E$.
\end{lemma}

\begin{proof}
Cela résulte directement de la définition.
\end{proof}

\noindent
Calculons maintenant le degré dans l'exemple particulier le plus important, celui
fourni par le Lemme \ref{lemma-field-extension-generated-by-one-element}, au moyen
des deux exemples suivants.

\begin{example}[Degré d'un corps de fonctions rationnelles]
\label{example-degree-rational-function-field}
Si $k$ est un corps quelconque, le corps des fonctions rationnelles $k(t)$ n'est
{\it pas} une extension finie. Par exemple, les éléments
$\left\{t^n, n \in \mathbf{Z}\right\}$ sont linéairement indépendants sur $k$.

\medskip\noindent
En fait, si $k$ est non dénombrable, alors $k(t)$ est de dimension {\it non dénombrable}
comme $k$-espace vectoriel. Pour le montrer, affirmons que la famille d'éléments
$\{1/(t- \alpha), \alpha \in k\} \subset k(t)$ est linéairement indépendante sur
$k$. Une relation non triviale entre eux conduirait à une contradiction : par
exemple, si l'on travaille sur $\mathbf{C}$, cela résulte du fait que
$\frac{1}{t-\alpha}$, considéré comme une fonction méromorphe sur
$\mathbf{C}$, possède un pôle en $\alpha$ et nulle part ailleurs.
Par conséquent, toute somme $\sum c_i \frac{1}{t - \alpha_i}$, où les $c_i \in k^*$
et les $\alpha_i \in k$ sont distincts, aurait un pôle en chacun des $\alpha_i$.
Elle ne pourrait notamment pas être nulle.

\medskip\noindent
Fait amusant, ceci donne une démonstration rapide du Nullstellensatz de Hilbert sur
les nombres complexes. Pour un résultat légèrement plus général, voir
Algèbre, Théorème \ref{algebra-theorem-uncountable-nullstellensatz}.
\end{example}

\begin{lemma}
\label{lemma-finite-finitely-generated}
Toute extension finie de corps est une extension de corps de type fini.
La réciproque est fausse.
\end{lemma}

\begin{proof}
Soit $F/E$ une extension finie de corps. Soient $\alpha_1, \ldots, \alpha_n$
une base de $F$ comme espace vectoriel sur $E$. Alors
$F = E(\alpha_1, \ldots, \alpha_n)$ ; ainsi $F/E$ est une extension de corps
de type fini. La réciproque est fausse d'après
l'Exemple \ref{example-degree-rational-function-field}.
\end{proof}

\begin{example}[Degré d'une extension algébrique simple]
\label{example-degree-simple-algebraic-extension}
Considérons une extension de corps monogène $E/k$ de la forme étudiée dans
l'Exemple \ref{example-monogenic-extension}.
Autrement dit, $E = k[t]/(P)$, où $P \in k[t]$ est un polynôme irréductible.
Alors le degré $[E : k]$ est simplement le degré $d = \deg(P)$ du
polynôme $P$. Écrivons en effet
\begin{equation}
\label{equation-P}
P = a_d t^d + a_{d - 1} t^{d - 1} + \ldots + a_0.
\end{equation}
avec $a_d \not = 0$. Les images de $1, t, \ldots, t^{d - 1}$ dans
$k[t]/(P)$ sont alors linéairement indépendantes sur $k$, car toute relation entre
elles serait de degré strictement inférieur à celui de $P$, alors que $P$ est
l'élément de plus petit degré de l'idéal $(P)$.

\medskip\noindent
Réciproquement, l'ensemble $S = \{1, t, \ldots, t^{d - 1}\}$ (ou, plus
précisément, l'ensemble de leurs images) engendre $k[t]/(P)$ comme espace vectoriel.
En effet, il résulte de (\ref{equation-P}) que $a_d t^d$ appartient au sous-espace engendré par $S$.
Comme $a_d$ est inversible, on voit que $t^d$ appartient au sous-espace engendré par $S$.
De même, la relation $t P(t) = 0$ montre que l'image de $t^{d + 1}$
appartient au sous-espace engendré par $\{1, t, \ldots, t^d\}$ --- donc, d'après ce qui précède,
au sous-espace engendré par $S$. Une récurrence ascendante montre alors
que l'image de $t^n$, pour $n \geq d$, appartient au sous-espace engendré par $S$.
\end{example}

\noindent
Ceci confirme, par exemple, que $[\mathbf{C}: \mathbf{R}] = 2$. Plus généralement,
si $k$ est un corps et $\alpha \in k$ n'est pas un carré, le polynôme irréductible
$x^2 - \alpha \in k[x]$ permet de construire une extension $k[x]/(x^2 - \alpha)$
de degré deux.
Nous noterons celle-ci $k(\sqrt{\alpha})$. De telles extensions seront dites
{\it quadratiques,} pour des raisons évidentes.

\medskip\noindent
Le fait fondamental concernant le degré est qu'il est {\it multiplicatif dans les tours.}

\begin{lemma}[Multiplicativité]
\label{lemma-multiplicativity-degrees}
Supposons donnée une tour de corps $F/E/k$. Alors
$$
[F:k] = [F:E][E:k]
$$
\end{lemma}

\begin{proof}
Soient $\alpha_1, \ldots, \alpha_n \in F$ une $E$-base de $F$, et
$\beta_1, \ldots, \beta_m \in E$ une $k$-base de $E$. L'affirmation est
que l'ensemble des produits
$\{\alpha_i \beta_j, 1 \leq i \leq n, 1 \leq j \leq m\}$
est une $k$-base de $F$. Vérifions d'abord qu'ils engendrent $F$ sur $k$.

\medskip\noindent
Par hypothèse, les $\{\alpha_i\}$ engendrent $F$ sur $E$. Ainsi, si
$f \in F$, il existe des $a_i \in E$ tels que
$$
f = \sum\nolimits_i a_i \alpha_i,
$$
et, pour chaque $i$, on peut écrire $a_i = \sum b_{ij} \beta_j$ avec des
$b_{ij} \in k$. En regroupant ces égalités, on obtient
$$
f = \sum\nolimits_{i,j} b_{ij} \alpha_i \beta_j,
$$
ce qui prouve que les $\{\alpha_i \beta_j\}$ engendrent $F$ sur $k$.

\medskip\noindent
Supposons maintenant qu'il existe une relation non triviale
$$
\sum\nolimits_{i,j} c_{ij} \alpha_i \beta_j = 0
$$
avec des $c_{ij} \in k$. On aurait alors
$$
\sum\nolimits_i \alpha_i \left( \sum\nolimits_j c_{ij} \beta_j \right) = 0,
$$
et les termes entre parenthèses appartiennent à $E$, puisque les $\beta_j$ y appartiennent.
L'indépendance $E$-linéaire des $\{\alpha_i\}$ montre que toutes les sommes intérieures sont nulles.
Puis l'indépendance $k$-linéaire des $\{\beta_j\}$ montre que tous les
$c_{ij}$ sont nuls.
\end{proof}

\noindent
Ouvrons une parenthèse pour donner une définition quelque peu tangente à notre propos.

\begin{definition}
\label{definition-number-field}
On dit qu'un corps $K$ est un {\it corps de nombres} s'il est de caractéristique
$0$ et si l'extension $K/\mathbf{Q}$ est finie.
\end{definition}

\noindent
Les corps de nombres sont les objets fondamentaux de la théorie algébrique des nombres. Nous verrons
plus loin que,
pour l'analogue des entiers $\mathbf{Z}$ dans un corps de nombres, une propriété
assez proche de la factorisation unique subsiste (bien qu'il n'y ait généralement
pas de factorisation unique au sens strict !).


\section{Extensions algébriques}
\label{section-algebraic-extensions}

\noindent
Une classe importante d'extensions est formée de celles dans lesquelles chaque élément engendre
une extension finie.

\begin{definition}
\label{definition-algebraic}
Considérons une extension de corps $F/E$. Un élément $\alpha \in F$ est dit
{\it algébrique} sur $E$ si $\alpha$ est racine d'un polynôme non nul
à coefficients dans $E$. Si tous les éléments de $F$ sont algébriques, on dit que $F$ est
une {\it extension algébrique} de $E$.
\end{definition}

\noindent
D'après le Lemme \ref{lemma-field-extension-generated-by-one-element}, la
sous-extension $E(\alpha)$ est isomorphe soit au corps des fonctions rationnelles
$E(t)$, soit à un anneau quotient $E[t]/(P)$, où $P \in E[t]$ est un
polynôme irréductible. Dans le second cas, $\alpha$ est algébrique sur
$E$ (en fait, la preuve du
Lemme \ref{lemma-field-extension-generated-by-one-element}
montre que l'on peut choisir $P$ de sorte que $\alpha$ soit une racine de $P$) ;
dans le premier cas, il ne l'est pas.

\begin{example}
\label{example-C-algebraic-over-R}
Le corps $\mathbf{C}$ est algébrique sur $\mathbf{R}$. En effet, si
$\alpha = a + ib$ dans $\mathbf{C}$, alors $\alpha^2 - 2a\alpha + a^2 + b^2 = 0$
est une équation polynomiale satisfaite par $\alpha$ sur $\mathbf{R}$.
\end{example}

\begin{example}
\label{example-compact-riemann-surface-is-finite-over-P1}
Soient $X$ une surface de Riemann compacte et
$f \in \mathbf{C}(X) - \mathbf{C}$ une fonction méromorphe non constante quelconque
sur $X$ (voir l'Exemple \ref{example-field-of-meromorphic-functions}). On sait alors
que $\mathbf{C}(X)$ est algébrique sur la sous-extension
$\mathbf{C}(f)$ engendrée par $f$. Nous ne le démontrerons pas.
\end{example}

\begin{lemma}
\label{lemma-algebraic-goes-up}
Soit $K/E/F$ une tour d'extensions de corps.
\begin{enumerate}
\item Si $\alpha \in K$ est algébrique sur $F$, alors $\alpha$ est algébrique
sur $E$.
\item Si $K$ est algébrique sur $F$, alors $K$ est algébrique sur $E$.
\end{enumerate}
\end{lemma}

\begin{proof}
Cela résulte immédiatement des définitions.
\end{proof}

\noindent
Montrons à présent qu'il existe un lien étroit entre la finitude et le caractère
algébrique.

\begin{lemma}
\label{lemma-finite-is-algebraic}
Toute extension finie est algébrique. En fait, une extension $E/k$ est algébrique
si et seulement si toute sous-extension $k(\alpha)/k$ engendrée par un
$\alpha \in E$ est finie.
\end{lemma}

\noindent
En général, il est tout à fait faux qu'une extension algébrique soit finie.

\begin{proof}
Soit $E/k$ finie, de degré $n$. Choisissons $\alpha \in E$. Les
éléments $1, \alpha, \ldots, \alpha^n$ sont linéairement
dépendants sur $k$, sans quoi on aurait nécessairement $[E : k] > n$. Une relation de
dépendance linéaire fournit alors le polynôme recherché que $\alpha$ doit annuler.

\medskip\noindent
Pour la dernière assertion, remarquons qu'une extension monogène $k(\alpha)/k$ est
finie si et seulement si $\alpha$ est algébrique sur $k$, d'après les
Exemples \ref{example-degree-rational-function-field} et
\ref{example-degree-simple-algebraic-extension}.
Ainsi, si $E/k$ est algébrique, chaque $k(\alpha)/k$, pour $\alpha \in E$, est une extension
finie, et réciproquement.
\end{proof}

\noindent
Nous pouvons extraire un lemme de la preuve précédente (en réalité, des
Exemples \ref{example-degree-rational-function-field} et
\ref{example-degree-simple-algebraic-extension}) :
une extension monogène est finie si et seulement si elle est algébrique.
Nous utiliserons cette observation dans le résultat suivant.

\begin{lemma}
\label{lemma-algebraic-finitely-generated}
\begin{slogan}
Une extension algébrique de type fini est finie.
\end{slogan}
Soient $k$ un corps et $\alpha_1, \alpha_2, \ldots, \alpha_n$ des éléments
d'un corps d'extension tels que chaque $\alpha_i$ soit algébrique sur $k$. Alors
l'extension $k(\alpha_1, \ldots, \alpha_n)/k$ est finie.
Autrement dit, une extension algébrique de type fini est finie.
\end{lemma}

\begin{proof}
En effet, chaque extension
$k(\alpha_{1}, \ldots, \alpha_{i+1})/k(\alpha_1, \ldots, \alpha_{i})$
est engendrée par un élément et algébrique, donc finie.
La multiplicativité du degré (Lemme \ref{lemma-multiplicativity-degrees})
donne le résultat.
\end{proof}

\noindent
L'ensemble des nombres complexes qui sont algébriques sur $\mathbf{Q}$ est simplement
appelé l'ensemble des {\it nombres algébriques.} Par exemple, $\sqrt{2}$ est algébrique,
$i$ est algébrique, mais $\pi$ ne l'est pas.
C'est un fait fondamental que les nombres algébriques forment un corps, bien qu'il ne soit pas
évident de le démontrer à partir de la définition selon laquelle un nombre est algébrique
exactement lorsqu'il satisfait une équation polynomiale non nulle à coefficients
rationnels (par exemple, en manipulant des équations polynomiales).

\begin{lemma}
\label{lemma-algebraic-elements}
Soit $E/k$ une extension de corps. Les éléments de $E$ algébriques sur $k$
forment une sous-extension de $E/k$.
\end{lemma}

\begin{proof}
Soient $\alpha, \beta \in E$ algébriques sur $k$. Alors $k(\alpha, \beta)/k$
est une extension finie d'après le Lemme \ref{lemma-algebraic-finitely-generated}.
Il s'ensuit que $k(\alpha + \beta) \subset k(\alpha, \beta)$ est une extension
finie, ce qui implique que $\alpha + \beta$ est algébrique d'après le
Lemme \ref{lemma-finite-is-algebraic}. Il en va de même pour la différence,
le produit et, lorsque $\beta \neq 0$, le quotient $\alpha/\beta$.
\end{proof}

\noindent
De nombreuses propriétés remarquables des extensions de corps, tout comme celles des anneaux, sont
préservées par les tours et les corps composés.

\begin{lemma}
\label{lemma-algebraic-permanence}
Soient $E/k$ et $F/E$ des extensions algébriques de corps. Alors $F/k$ est une
extension algébrique de corps.
\end{lemma}

\begin{proof}
Choisissons $\alpha \in F$. Alors $\alpha$ est algébrique sur $E$.
L'observation essentielle est que $\alpha$ est algébrique sur une
sous-extension de type fini de $k$.
Autrement dit, il existe un ensemble fini $S \subset E$ tel que $\alpha $ soit algébrique
sur $k(S)$ : cela est clair, puisque le fait d'être algébrique signifie qu'il existe un
polynôme de $E[x]$ annulé par $\alpha$, et l'on peut prendre pour $S$ l'ensemble des
coefficients de ce polynôme. Il en résulte que $\alpha$ est algébrique sur
$k(S)$. En particulier, l'extension $k(S, \alpha)/ k(S)$ est finie.
Comme $S$ est un ensemble fini et que $k(S)/k$ est algébrique, le
Lemme \ref{lemma-algebraic-finitely-generated} montre que
$k(S)/k$ est finie. En utilisant la multiplicativité
(Lemme \ref{lemma-multiplicativity-degrees}),
on trouve que $k(S,\alpha)/k$ est finie ; ainsi $\alpha$ est algébrique sur $k$.
\end{proof}

\noindent
La méthode de la preuve précédente --- le fait que la propriété d'être algébrique
sur $E$ {\it descendait} à une sous-extension de type fini
de $E$ --- est une idée récurrente en algèbre.
Elle permet souvent, par exemple, de ramener des questions générales d'algèbre commutative
au cas noethérien.

\begin{lemma}
\label{lemma-size-algebraic-extension}
Soit $E/F$ une extension algébrique de corps. Alors le cardinal $|E|$
de $E$ est au plus $\max(\aleph_0, |F|)$.
\end{lemma}

\begin{proof}
Soit $S$ l'ensemble des polynômes non constants à coefficients dans $F$.
Pour tout $P \in S$, l'ensemble des racines
$r(P, E) = \{\alpha \in E \mid P(\alpha) = 0\}$
est fini (les détails sont omis). De plus, le fait que $E$ soit algébrique
sur $F$ implique que $E = \bigcup_{P \in S} r(P, E)$.
Il est clair que le cardinal de $S$ est majoré par $\max(\aleph_0, |F|)$,
car cet ensemble est une réunion dénombrable de produits finis de copies de $F$.
Il en va donc de même pour $E$.
\end{proof}

\begin{lemma}
\label{lemma-subalgebra-algebraic-extension-field}
Soit $E/F$ une extension finie, ou plus généralement algébrique, de corps.
Tout sous-anneau $F \subset R \subset E$ est un corps.
\end{lemma}

\begin{proof}
Soit $\alpha \in R$ non nul. Alors $1, \alpha, \alpha^2, \ldots$
appartiennent à $R$. D'après le Lemme \ref{lemma-finite-is-algebraic},
il existe une relation non triviale
$a_0 + a_1 \alpha + \ldots + a_d \alpha^d = 0$, avec $a_i \in F$.
Nous pouvons supposer $a_0 \not = 0$, car sinon nous pouvons diviser la relation
par $\alpha$ afin de diminuer $d$. On obtient alors
$$
a_0 = \alpha (- a_1  - \ldots - a_d \alpha^{d - 1})
$$
ce qui prouve que l'inverse de $\alpha$ est l'élément
$a_0^{-1} (- a_1  - \ldots - a_d \alpha^{d - 1})$
de $R$.
\end{proof}

\begin{lemma}
\label{lemma-algebraic-extension-self-map}
Soit $E/F$ une extension algébrique de corps. Tout morphisme de $F$-algèbres
$f : E \to E$ est un automorphisme.
\end{lemma}

\begin{proof}
Si $E/F$ est finie, alors $f : E \to E$ est une application $F$-linéaire
injective (Lemme \ref{lemma-field-maps-injective})
entre espaces vectoriels de dimension finie, donc est bijective.
Dans le cas général, on voit encore que $f$ est injective.
Soit $\alpha \in E$, et soit $P \in F[x]$ un
polynôme tel que $P(\alpha) = 0$.
Soit $E' \subset E$ le sous-corps de $E$ engendré
par les racines $\alpha = \alpha_1, \ldots, \alpha_n$ de $P$ dans $E$.
Alors $E'$ est fini sur $F$ d'après le Lemme \ref{lemma-algebraic-finitely-generated}.
Comme $f$ préserve l'ensemble des racines, on a
$f|_{E'} : E' \to E'$. Ainsi $f|_{E'}$ est un isomorphisme
d'après la première partie de la preuve, et l'on conclut que $\alpha$
appartient à l'image de $f$.
\end{proof}








\section{Polynômes minimaux}
\label{section-minimal-polynomials}

\noindent
Soient $E/k$ une extension de corps et $\alpha \in E$ algébrique sur $k$.
Alors $\alpha$ satisfait une équation polynomiale (non triviale) dans $k[x]$.
Considérons l'ensemble des polynômes $P \in k[x]$ tels que $P(\alpha) = 0$ ; par
hypothèse, cet ensemble ne contient pas seulement le polynôme nul.
On vérifie facilement que cet ensemble est un {\it idéal.} C'est en effet le noyau
du morphisme
$$
k[x] \to E, \quad x \mapsto \alpha
$$
Puisque $k[x]$ est un anneau principal, il existe un {\it générateur} $P \in k[x]$ de cet
idéal. Si l'on suppose, sans perte de généralité, que $P$ est unitaire, alors $P$ est
déterminé de manière unique.

\begin{definition}
\label{definition-minimal-polynomial}
Le polynôme $P$ ci-dessus est appelé le {\it polynôme minimal}
de $\alpha$ sur $k$.
\end{definition}

\noindent
Le polynôme minimal possède la caractérisation suivante : c'est le polynôme unitaire,
de plus petit degré, qui annule $\alpha$. Tout multiple non constant de $P$
sera de degré plus grand, et seuls les multiples de $P$ peuvent
annuler $\alpha$. Cela explique le qualificatif {\it minimal}.

\medskip\noindent
Le polynôme minimal est manifestement {\it irréductible}. Cela équivaut à affirmer
que l'idéal de $k[x]$ formé des polynômes qui annulent
$\alpha$ est premier. Ce fait résulte de ce que le morphisme
$k[x] \to E, x \mapsto \alpha$ prend ses valeurs dans un anneau intègre (et même dans un corps), de sorte que son
noyau est un idéal premier.

\begin{lemma}
\label{lemma-degree-minimal-polynomial}
Le degré du polynôme minimal est $[k(\alpha) : k]$.
\end{lemma}

\begin{proof}
Il s'agit simplement d'une reformulation de l'argument donné dans le
Lemme \ref{lemma-field-extension-generated-by-one-element} : l'observation
est que, si $P$ est le polynôme minimal de $\alpha$, le morphisme
$$
k[x]/(P) \to k(\alpha), \quad x \mapsto \alpha
$$
est un isomorphisme comme dans la preuve précitée, et nous avons calculé le
degré d'une telle extension (voir
l'Exemple \ref{example-degree-simple-algebraic-extension}).
\end{proof}

\noindent
L'observation de la preuve ci-dessus est donc que, si $\alpha \in E$ est algébrique,
alors $k(\alpha) \subset E$ est isomorphe à $k[x]/(P)$.


\section{Clôture algébrique}
\label{section-algebraic-closure}

\noindent
Le ``théorème fondamental de l'algèbre'' affirme que $\mathbf{C}$ est
algébriquement clos. Une belle démonstration de ce résultat utilise
le théorème de Liouville en analyse complexe ; nous en donnerons une autre
(voir le Lemme \ref{lemma-C-algebraically-closed}).

\begin{definition}
\label{definition-algebraically-closed}
On dit qu'un corps $F$ est {\it algébriquement clos} si toute extension algébrique
$E/F$ est triviale, c'est-à-dire si $E = F$.
\end{definition}

\noindent
Ce n'est peut-être pas la définition retenue dans tous les ouvrages. Le lemme suivant
la compare à l'autre définition usuelle.

\begin{lemma}
\label{lemma-algebraically-closed}
Soit $F$ un corps. Les assertions suivantes sont équivalentes :
\begin{enumerate}
\item $F$ est algébriquement clos,
\item tout polynôme irréductible sur $F$ est de degré un,
\item tout polynôme non constant sur $F$ possède une racine,
\item tout polynôme non constant sur $F$ est un produit de facteurs de degré un.
\end{enumerate}
\end{lemma}

\begin{proof}
Si $F$ est algébriquement clos, alors tout polynôme irréductible est de degré un.
En effet, s'il existe un polynôme irréductible de degré $> 1$, celui-ci
engendre une extension de corps finie non triviale (donc algébrique) ; voir
l'Exemple \ref{example-degree-simple-algebraic-extension}.
Ainsi (1) implique (2). Si tout polynôme irréductible
est de degré un, alors tout polynôme irréductible possède une racine, et par conséquent tout
polynôme non constant possède une racine. Ainsi (2) implique (3).

\medskip\noindent
Supposons que tout polynôme non constant possède une racine. Soit $P \in F[x]$
non constant. Si $P(\alpha) = 0$ avec $\alpha \in F$, on voit alors
que $P = (x - \alpha)Q$ pour un certain $Q \in F[x]$ (par division euclidienne).
On peut donc raisonner par récurrence sur le degré pour écrire tout polynôme non constant
sous la forme d'un produit $c \prod (x - \alpha_i)$.

\medskip\noindent
Supposons enfin que tout polynôme non constant sur $F$ soit un produit de
facteurs de degré un. Soit $E/F$ une extension algébrique. Toutes les sous-extensions
simples $F(\alpha)/F$ de $E$ sont alors nécessairement triviales (car, par
hypothèse, les seuls polynômes irréductibles sont de degré un). Ainsi $E = F$.
Nous voyons que (4) implique (1), ce qui achève la preuve.
\end{proof}

\noindent
Nous voulons maintenant définir une extension algébrique ``universelle'' d'un corps.
En réalité, il convient d'être prudent : la clôture algébrique n'est {\it pas} un
objet universel. En effet, la clôture algébrique n'est pas unique à
{\it unique} isomorphisme près : elle n'est unique qu'à isomorphisme près. Elle
n'en sera pas moins très utile, bien que non fonctorielle.

\begin{definition}
\label{definition-algebraic-closure}
Soit $F$ un corps. Une {\it clôture algébrique} de $F$ est un corps
$\overline{F}$ contenant $F$ tel que :
\begin{enumerate}
\item $\overline{F}$ est algébrique sur $F$.
\item $\overline{F}$ est algébriquement clos.
\end{enumerate}
\end{definition}

\noindent
Si $F$ est algébriquement clos, alors $F$ est sa propre clôture algébrique.
Démontrons maintenant le résultat fondamental d'existence.

\begin{theorem}
\label{theorem-existence-algebraic-closure}
Tout corps possède une clôture algébrique.
\end{theorem}

\noindent
La preuve sera pour l'essentiel sans rapport avec le reste du chapitre. Nous voudrons toutefois
savoir qu'il est {\it possible} de plonger un corps dans un
corps algébriquement clos, et nous supposerons souvent ce plongement effectué.

\begin{proof}
Soit $F$ un corps. D'après le Lemme \ref{lemma-size-algebraic-extension}, le
cardinal d'une extension algébrique de $F$ est majoré par
$\max(\aleph_0, |F|)$. Choisissons un ensemble $S$ contenant $F$ tel que
$|S| > \max(\aleph_0, |F|)$. Considérons les triplets
$(E, \sigma_E, \mu_E)$ où
\begin{enumerate}
\item $E$ est un ensemble tel que $F \subset E \subset S$, et
\item $\sigma_E : E \times E \to E$ et $\mu_E : E \times E \to E$
sont des applications d'ensembles telles que $(E, \sigma_E, \mu_E)$ définisse la structure
d'une extension de corps de $F$ (en particulier, $\sigma_E(a, b) = a +_F b$
pour $a, b \in F$, et de même pour $\mu_E$), et
\item $E/F$ est une extension algébrique de corps.
\end{enumerate}
La collection de tous les triplets $(E, \sigma_E, \mu_E)$ forme un ensemble $I$.
Pour $i \in I$, nous noterons $E_i = (E_i, \sigma_i, \mu_i)$ l'extension de corps
de $F$ correspondante. Nous définissons un ordre partiel sur
$I$ en déclarant que $i \leq i'$ si et seulement si $E_i \subset E_{i'}$
(cela a un sens puisque $E_i$ et $E_{i'}$ sont des sous-ensembles du même ensemble $S$)
et si $\sigma_i = \sigma_{i'}|_{E_i \times E_i}$ et
$\mu_i = \mu_{i'}|_{E_i \times E_i}$ ; autrement dit, $E_{i'}$ est une extension de corps
de $E_i$.

\medskip\noindent
Soit $T \subset I$ un sous-ensemble totalement ordonné. Il est alors clair que
$E_T = \bigcup_{i \in T} E_i$, muni des applications induites $\sigma_T = \bigcup \sigma_i$
et $\mu_T = \bigcup \mu_i$, est un autre élément de $I$. Autrement dit,
tout sous-ensemble totalement ordonné de $I$ possède un majorant dans $I$. D'après le lemme de Zorn,
il existe un élément maximal $(E, \sigma_E, \mu_E)$ de $I$. Nous affirmons que
$E$ est une clôture algébrique. Puisque, par définition de $I$, l'extension
$E/F$ est algébrique, il suffit de montrer que $E$ est algébriquement clos.

\medskip\noindent
Pour le voir, raisonnons par l'absurde. Supposons donc que $E$ ne soit pas
algébriquement clos. Il existe alors un polynôme irréductible
$P$ sur $E$ de degré $> 1$ ; voir le Lemme \ref{lemma-algebraically-closed}.
D'après le Lemme \ref{lemma-finite-is-algebraic}, on obtient une extension finie
non triviale $E' = E[x]/(P)$. Remarquons que $E'/F$ est algébrique d'après le
Lemme \ref{lemma-algebraic-permanence}.
Le cardinal de $E'$ est donc $\leq \max(\aleph_0, |F|)$.
La théorie élémentaire des ensembles permet de prolonger l'injection donnée
$E \subset S$ en une injection $E' \to S$. Autrement dit, nous pouvons
considérer $E'$ comme un élément de notre ensemble $I$, ce qui contredit la
maximalité de $E$. Cette contradiction achève la preuve.
\end{proof}

\begin{lemma}
\label{lemma-map-into-algebraic-closure}
Soit $F$ un corps. Soit $\overline{F}$ une clôture algébrique de $F$.
Soit $M/F$ une extension algébrique. Il existe alors un morphisme
d'extensions de $F$, $M \to \overline{F}$.
\end{lemma}

\begin{proof}
Considérons l'ensemble $I$ des couples $(E, \varphi)$, où $F \subset E \subset M$
est une sous-extension et $\varphi : E \to \overline{F}$ un morphisme
d'extensions de $F$. Munissons $I$ de l'ordre partiel défini en déclarant que
$(E, \varphi) \leq (E', \varphi')$ si et seulement si $E \subset E'$ et
$\varphi'|_E = \varphi$. Si $T = \{(E_t,  \varphi_t)\} \subset I$
est un sous-ensemble totalement ordonné, alors
$\bigcup \varphi_t : \bigcup E_t \to \overline{F}$ est un élément de $I$.
Ainsi, tout sous-ensemble totalement ordonné de $I$ possède un majorant.
Le lemme de Zorn assure l'existence d'un élément maximal $(E, \varphi)$ de $I$.
Nous affirmons que $E = M$, ce qui achèvera la preuve. Sinon,
choisissons $\alpha \in M$, $\alpha \not \in E$. Alors $\alpha$ est algébrique
sur $E$ ; voir le Lemme \ref{lemma-algebraic-goes-up}.
Soit $P$ le polynôme minimal de $\alpha$ sur $E$.
Soit $P^\varphi$ l'image de $P$ par $\varphi$ dans $\overline{F}[x]$.
Comme $\overline{F}$ est algébriquement clos, il existe une racine $\beta$
de $P^\varphi$ dans $\overline{F}$. On peut alors prolonger $\varphi$ en
$\varphi' : E(\alpha) = E[x]/(P) \to \overline{F}$ en envoyant
$x$ sur $\beta$. Cela contredit la maximalité de $(E, \varphi)$,
comme voulu.
\end{proof}

\begin{lemma}
\label{lemma-algebraic-closures-isomorphic}
Deux clôtures algébriques quelconques d'un corps sont isomorphes.
\end{lemma}

\begin{proof}
Soit $F$ un corps. Si $M$ et $\overline{F}$ sont des clôtures algébriques de
$F$, il existe alors un morphisme d'extensions de $F$
$\varphi : M \to \overline{F}$ d'après le
Lemme \ref{lemma-map-into-algebraic-closure}.
Or l'image $\varphi(M)$ est algébriquement close.
D'autre part, l'extension $\varphi(M) \subset \overline{F}$
est algébrique d'après le Lemme \ref{lemma-algebraic-goes-up}.
Ainsi $\varphi(M) = \overline{F}$.
\end{proof}





\section{Polynômes premiers entre eux}
\label{section-relatively-prime}

\noindent
Soit $K$ un corps algébriquement clos. L'anneau $K[x]$ possède alors une structure
d'idéaux très simple, comme nous l'avons vu au Lemme \ref{lemma-algebraically-closed}.
En particulier, tout polynôme $P \in K[x]$ peut s'écrire
$$
P = c(x - \alpha_1) \ldots (x - \alpha_n),
$$
où $c$ est le coefficient dominant et où les $\alpha_1, \ldots, \alpha_n \in K$
sont les racines de $P$ (comptées avec multiplicité). Manifestement, les seuls polynômes
irréductibles de $K[x]$ sont les polynômes de degré un $c(x - \alpha)$,
$c, \alpha \in K$ (avec $c \neq 0$).

\begin{definition}
\label{definition-relatively-prime}
Si $k$ est un corps quelconque, nous disons que deux polynômes de $k[x]$ sont
{\it premiers entre eux} s'ils engendrent l'idéal unité de $k[x]$.
\end{definition}

\noindent
Poursuivons la discussion précédente. Si $K$ est un corps algébriquement clos,
deux polynômes de $K[x]$ sont premiers entre eux si et seulement s'ils n'ont aucune
racine commune. Cela résulte du fait que les idéaux maximaux de $K[x]$ sont de la forme
$(x - \alpha)$, $\alpha \in K$. Ainsi, si $F, G \in K[x]$ n'ont aucune racine commune,
alors $(F, G)$ ne peut être contenu dans aucun $(x - \alpha)$ (car ils auraient alors
une racine commune en $\alpha$).

\medskip\noindent
Si $k$ n'est {\it pas} algébriquement clos, ceci fournit encore des
informations permettant de déterminer quand deux polynômes de $k[x]$ engendrent l'idéal unité.

\begin{lemma}
\label{lemma-relatively-prime-polynomials}
Deux polynômes de $k[x]$ sont premiers entre eux si et seulement s'ils
n'ont aucune racine commune dans une clôture algébrique $\overline{k}$ de $k$.
\end{lemma}

\begin{proof}
L'affirmation est que deux polynômes quelconques $P, Q$ engendrent $(1)$ dans $k[x]$ si et
seulement s'ils engendrent $(1)$ dans $\overline{k}[x]$. C'est une question
d'algèbre linéaire : un système d'équations linéaires à coefficients dans $k$ possède
une solution si et seulement s'il possède une solution dans une extension quelconque de $k$.
On peut par conséquent se ramener au cas d'un corps algébriquement clos ; dans
ce cas, le résultat découle clairement de ce qui précède.
\end{proof}





\section{Extensions algébriques séparables}
\label{section-separable-extensions}

\noindent
En caractéristique $p$, un phénomène curieux se produit pour les polynômes irréductibles
sur un corps. Le lemme suivant l'explique.

\begin{lemma}
\label{lemma-irreducible-polynomials}
Soit $F$ un corps. Soit $P \in F[x]$ un polynôme irréductible sur $F$.
Soit $P' = \text{d}P/\text{d}x$ le polynôme dérivé de $P$ par rapport
à $x$. L'un des deux cas suivants se produit :
\begin{enumerate}
\item $P$ et $P'$ sont premiers entre eux, ou
\item $P'$ est le polynôme nul.
\end{enumerate}
Le second cas ne peut se produire que si $F$ est de caractéristique $p > 0$.
Dans ce cas, $P(x) = Q(x^q)$, où $q = p^f$ est une puissance de $p$ et
$Q \in F[x]$ un polynôme irréductible tel que $Q$ et $Q'$ soient
premiers entre eux.
\end{lemma}

\begin{proof}
Remarquons que $P'$ est de degré $< \deg(P)$. Par conséquent, si $P$ et $P'$ ne sont pas premiers
entre eux, alors $(P, P') = (R)$, où $R$ est un polynôme de degré $< \deg(P)$,
ce qui contredit l'irréductibilité de $P$. Cela établit la dichotomie
entre (1) et (2).

\medskip\noindent
Supposons que nous soyons dans le cas (2) et que $P = a_d x^d + \ldots + a_0$. Alors
$P' = da_d x^{d - 1} + \ldots + a_1$. En caractéristique $0$, on voit
que cela impose $a_d, \ldots, a_1 = 0$, ce qui signifierait que $P$ est constant,
une contradiction. On conclut donc que la caractéristique $p$ est positive.
Dans ce cas, la condition $P' = 0$ impose $a_i = 0$ chaque fois que $p$ ne
divise pas $i$.
Autrement dit, $P(x) = P_1(x^p)$ pour un certain polynôme non constant $P_1$.
Manifestement, $P_1$ est lui aussi irréductible. Par récurrence sur le degré, on
voit que $P_1(x) = Q(x^q)$ comme dans l'énoncé du lemme ; ainsi
$P(x) = Q(x^{pq})$, et le lemme est démontré.
\end{proof}

\begin{definition}
\label{definition-separable}
Soit $F$ un corps. Soit $K/F$ une extension de corps.
\begin{enumerate}
\item On dit qu'un polynôme irréductible $P$ sur $F$ est {\it séparable}
s'il est premier avec son polynôme dérivé.
\item Étant donné $\alpha \in K$ algébrique sur $F$, on dit que $\alpha$ est
{\it séparable} sur $F$ si son polynôme minimal est séparable sur $F$.
\item Si $K$ est une extension algébrique de $F$, on dit que $K$ est
{\it séparable}\footnote{Pour les extensions non algébriques,
cette définition n'a pas de sens et n'est pas la bonne. Nous renvoyons
le lecteur à Algèbre, Sections \ref{algebra-section-separability} et
\ref{algebra-section-separability-continued}.}
sur $F$ si tout élément de $K$ est séparable sur $F$.
\end{enumerate}
\end{definition}

\noindent
D'après le Lemme \ref{lemma-irreducible-polynomials}, en caractéristique $0$, tout
polynôme irréductible est séparable, tout élément algébrique d'une extension
est séparable et toute extension algébrique est séparable.

\begin{lemma}
\label{lemma-separable-goes-up}
Soit $K/E/F$ une tour d'extensions algébriques de corps.
\begin{enumerate}
\item Si $\alpha \in K$ est séparable sur $F$, alors $\alpha$ est séparable
sur $E$.
\item Si $K$ est séparable sur $F$, alors $K$ est séparable sur $E$.
\end{enumerate}
\end{lemma}

\begin{proof}
Nous utiliserons le Lemme \ref{lemma-irreducible-polynomials} sans le mentionner davantage.
Soit $P$ le polynôme minimal de $\alpha$ sur $F$.
Soit $Q$ le polynôme minimal de $\alpha$ sur $E$.
Alors $Q$ divise $P$ dans l'anneau de polynômes $E[x]$, disons $P = QR$.
Alors $P' = Q'R + QR'$. Ainsi, si $Q' = 0$, alors $Q$ divise $P$ et $P'$,
donc $P' = 0$ d'après le lemme. Ceci démontre (1). L'assertion (2)
découle immédiatement de (1) et des définitions.
\end{proof}

\begin{lemma}
\label{lemma-recognize-separable}
Soit $F$ un corps. Un polynôme irréductible $P$ sur $F$
est séparable si et seulement si $P$ possède des racines deux à deux distinctes dans une
clôture algébrique de $F$.
\end{lemma}

\begin{proof}
Supposons que $\alpha \in \overline{F}$ soit une racine à la fois de $P$ et de $P'$.
Alors $P = (x - \alpha)Q$ pour un certain polynôme $Q$. En dérivant,
nous obtenons $P' = Q + (x - \alpha)Q'$. Ainsi $\alpha$ est une racine de $Q$.
On voit donc que si $P$ et $P'$ ont une racine commune, alors $P$
ne possède pas de racines deux à deux distinctes. Réciproquement, si $P$ possède
une racine multiple, c'est-à-dire si $(x - \alpha)^2$ divise $P$, alors $\alpha$
est une racine à la fois de $P$ et de $P'$. Conjointement avec le
Lemme \ref{lemma-relatively-prime-polynomials}, ceci démontre le lemme.
\end{proof}

\begin{lemma}
\label{lemma-nr-roots-unchanged}
Soit $F$ un corps et soit $\overline{F}$ une clôture algébrique de $F$.
Soit $p > 0$ la caractéristique de $F$. Soit $P$ un polynôme
sur $F$. Alors les ensembles des racines de $P$ et de $P(x^p)$ dans $\overline{F}$
ont le même cardinal (sans compter les multiplicités).
\end{lemma}

\begin{proof}
Manifestement, $\alpha$ est une racine de $P(x^p)$ si et seulement si $\alpha^p$ est une
racine de $P$. Autrement dit, les racines de $P(x^p)$ sont les racines de
$x^p - \beta$, où $\beta$ est une racine de $P$. Il suffit donc de montrer
que l'application $\overline{F} \to \overline{F}$, $\alpha \mapsto \alpha^p$
est bijective. Elle est surjective, puisque $\overline{F}$ est algébriquement clos,
ce qui signifie que tout élément possède une racine $p$-ième. Elle est injective car
$\alpha^p = \beta^p$ implique $(\alpha - \beta)^p = 0$, puisque
la caractéristique est $p$. Et bien sûr, dans un corps, $x^p = 0$ implique
$x = 0$.
\end{proof}

\noindent
Soit $F$ un corps et soit $P$ un polynôme irréductible sur $F$.
Nous savons alors que $P = Q(x^q)$ pour un certain polynôme irréductible séparable $Q$
(Lemme \ref{lemma-irreducible-polynomials}), où $q$ est une puissance de
la caractéristique $p$ (et, si la caractéristique est nulle, alors
$q = 1$\footnote{Une bonne convention pour ce chapitre consiste à poser $0^0 = 1$.}
et $Q = P$). D'après le Lemme \ref{lemma-nr-roots-unchanged}, le nombre de
racines de $P$ et de $Q$ dans toute clôture algébrique de $F$ est le même.
D'après le Lemme \ref{lemma-recognize-separable}, ce nombre est égal au degré
de $Q$.

\begin{definition}
\label{definition-separable-degree}
Soit $F$ un corps. Soit $P$ un polynôme irréductible sur $F$.
Le {\it degré séparable} de $P$ est le cardinal de
l'ensemble des racines de $P$ dans toute clôture algébrique de $F$ (voir la discussion
ci-dessus). Notation $\deg_s(P)$.
\end{definition}

\noindent
Le degré séparable de $P$ divise toujours son degré, et le quotient
est une puissance de la caractéristique. Si la caractéristique est nulle, alors
$\deg_s(P) = \deg(P)$.

\begin{situation}
\label{situation-finitely-generated}
Ici, $F$ est un corps et $K/F$ est une extension finie engendrée par des éléments
$\alpha_1, \ldots, \alpha_n \in K$. Nous posons $K_0 = F$ et
$$
K_i = F(\alpha_1, \ldots, \alpha_i)
$$
pour obtenir une tour d'extensions finies
$K = K_n / K_{n - 1} / \ldots / K_0 = F$.
Notons $P_i$ le polynôme minimal de $\alpha_i$ sur $K_{i - 1}$.
Enfin, fixons une clôture algébrique $\overline{F}$ de $F$.
\end{situation}

\noindent
Soient $F$, $K$, les $\alpha_i$ et $\overline{F}$ comme dans la
Situation \ref{situation-finitely-generated}.
Supposons que $\varphi : K \to \overline{F}$ soit un morphisme d'extensions
de $F$. Nous obtenons alors des applications $\varphi_i : K_i \to \overline{F}$.
L'image de $P_i \in K_{i - 1}[x]$ par $\varphi_{i - 1}$
est un polynôme $P_i^{\varphi_{i - 1}} \in \overline{F}[x]$
et $\varphi(\alpha_i)$ est une racine de ce polynôme.

\begin{lemma}
\label{lemma-count-embeddings}
Dans la Situation \ref{situation-finitely-generated}, la correspondance
$$
\Mor_F(K, \overline{F})
\longrightarrow
\{(\beta_1, \ldots, \beta_n)\text{ comme ci-dessous}\},
\quad
\varphi \longmapsto (\varphi(\alpha_1), \ldots, \varphi(\alpha_n))
$$
est une bijection. Ici, le membre de droite est l'ensemble des $n$-uplets
$(\beta_1, \ldots, \beta_n)$ d'éléments de $\overline{F}$ possédant
la propriété suivante :
\begin{enumerate}
\item $\beta_1 \in \overline{F}$ est une racine de $P_1$ ;
soit $\varphi_1 : K_1 \to \overline{F}$ l'homomorphisme
de $F$-algèbres qui envoie $\alpha_1$ sur $\beta_1$,
\item $\beta_2 \in \overline{F}$ est une racine de $P_2^{\varphi_1}$ ;
soit $\varphi_2 : K_2 \to \overline{F}$ l'homomorphisme
qui prolonge $\varphi_1$ et envoie $\alpha_2$ sur $\beta_2$,
\item et ainsi de suite jusqu'au dernier cas,
\item $\beta_n \in \overline{F}$ est une racine de $P_n^{\varphi_{n - 1}}$.
\end{enumerate}
À chaque étape, l'homomorphisme $\varphi_i$ existe et est unique, car
$K_i = K_{i - 1}[x]/(P_i)$ et $\beta_i$ est une racine de $P_i^{\varphi_{i - 1}}$.
\end{lemma}

\begin{proof}
L'application de gauche à droite a été décrite avant le lemme.
Soit $(\beta_1, \ldots, \beta_n)$ un élément du membre de droite.
Nous définissons alors $\varphi : K = K_n \to \overline{F}$ comme l'unique homomorphisme
qui prolonge $\varphi_{n - 1}$ et envoie $\alpha_n$ sur $\beta_n$.
L'unicité implique que les deux constructions sont inverses l'une de l'autre.
\end{proof}

\begin{lemma}
\label{lemma-count-embeddings-explicitly}
Dans la Situation \ref{situation-finitely-generated}, nous avons
$|\Mor_F(K, \overline{F})| = \prod_{i = 1}^n \deg_s(P_i)$.
\end{lemma}

\begin{proof}
Ceci découle immédiatement du Lemme \ref{lemma-count-embeddings}.
Remarquons qu'un ingrédient essentiel que nous utilisons ici implicitement est le
caractère bien défini du degré séparable d'un polynôme irréductible,
qui a été constaté juste avant la
Définition \ref{definition-separable-degree}.
\end{proof}

\noindent
Nous utilisons maintenant le résultat précédent pour caractériser les extensions de corps séparables.

\begin{lemma}
\label{lemma-separably-generated-separable}
Hypothèses et notations comme dans la Situation \ref{situation-finitely-generated}.
Si chaque $P_i$ est séparable, c'est-à-dire si chaque $\alpha_i$ est séparable sur
$K_{i - 1}$, alors
$$
|\Mor_F(K, \overline{F})| = [K : F]
$$
et l'extension de corps $K/F$ est séparable. Si l'un des $\alpha_i$ n'est
pas séparable sur $K_{i - 1}$, alors
$|\Mor_F(K, \overline{F})| < [K : F]$.
\end{lemma}

\begin{proof}
Si $\alpha_i$ est séparable sur $K_{i - 1}$, alors
$\deg_s(P_i) = \deg(P_i) = [K_i : K_{i - 1}]$
(la dernière égalité résultant du Lemme \ref{lemma-degree-minimal-polynomial}).
Par multiplicativité (Lemme \ref{lemma-multiplicativity-degrees}), nous avons
$$
[K : F] = \prod [K_i : K_{i - 1}] = \prod \deg(P_i) =
\prod \deg_s(P_i) = |\Mor_F(K, \overline{F})|
$$
où la dernière égalité est le Lemme \ref{lemma-count-embeddings-explicitly}.
Par exactement le même raisonnement, nous obtenons l'inégalité stricte
$|\Mor_F(K, \overline{F})| < [K : F]$ si l'un des $\alpha_i$ n'est
pas séparable sur $K_{i - 1}$.

\medskip\noindent
Supposons enfin, à nouveau, que chaque $\alpha_i$ soit séparable sur $K_{i - 1}$.
Nous allons montrer que $K/F$ est séparable.
Soit $\gamma = \gamma_1 \in K$ quelconque. Nous pouvons alors trouver des éléments supplémentaires
$\gamma_2, \ldots, \gamma_m$ tels que
$K = F(\gamma_1, \ldots, \gamma_m)$ (nous pourrions par exemple prendre
$\gamma_2 = \alpha_1, \ldots, \gamma_{n + 1} = \alpha_n$).
La dernière partie du lemme (déjà démontrée ci-dessus) montre alors
que si $\gamma$ n'était pas séparable sur $F$, nous aurions l'inégalité
stricte $|\Mor_F(K, \overline{F})| < [K : F]$,
en contradiction avec la toute première partie du lemme (elle aussi déjà démontrée
ci-dessus).
\end{proof}

\begin{lemma}
\label{lemma-separable-equality}
Soit $K/F$ une extension finie de corps. Soit $\overline{F}$ une
clôture algébrique de $F$. Nous avons alors
$$
|\Mor_F(K, \overline{F})| \leq [K : F]
$$
avec égalité si et seulement si $K$ est séparable sur $F$.
\end{lemma}

\begin{proof}
C'est un corollaire du Lemme \ref{lemma-separably-generated-separable}.
En effet, puisque $K/F$ est finie, nous pouvons trouver un nombre fini d'éléments
$\alpha_1, \ldots, \alpha_n \in K$ qui engendrent $K$ sur $F$ (nous pouvons par exemple
choisir les $\alpha_i$ de manière à former une base de $K$ sur $F$).
Si $K/F$ est séparable, alors chaque $\alpha_i$ est séparable sur
$F(\alpha_1, \ldots, \alpha_{i - 1})$ d'après le Lemme \ref{lemma-separable-goes-up},
et le Lemme \ref{lemma-separably-generated-separable} donne l'égalité.
Réciproquement, si nous avons l'égalité, alors, quelle que soit la manière de choisir
$\alpha_1, \ldots, \alpha_n$, le Lemme \ref{lemma-separably-generated-separable}
montre que $\alpha_1$ est séparable sur $F$. Comme nous pouvons
commencer la suite par un élément arbitraire de $K$, il s'ensuit
que $K$ est séparable sur $F$.
\end{proof}

\begin{lemma}
\label{lemma-separable-permanence}
Soient $E/k$ et $F/E$ des extensions algébriques séparables de corps. Alors $F/k$
est une extension séparable de corps.
\end{lemma}

\begin{proof}
Choisissons $\alpha \in F$. Alors $\alpha$ est algébrique séparable sur $E$.
Soit $P = x^d + \sum_{i < d} a_i x^i$ le polynôme minimal de
$\alpha$ sur $E$. Chaque $a_i$ est algébrique séparable sur $k$.
Considérons la tour de corps
$$
k \subset k(a_0) \subset k(a_0, a_1) \subset \ldots \subset
k(a_0, \ldots, a_{d - 1}) \subset k(a_0, \ldots, a_{d - 1}, \alpha)
$$
Puisque $a_i$ est algébrique séparable sur $k$, il est algébrique séparable
sur $k(a_0, \ldots, a_{i - 1})$ d'après le Lemme \ref{lemma-separable-goes-up}.
Enfin, $\alpha$ est algébrique séparable sur $k(a_0, \ldots, a_{d - 1})$
car c'est une racine de $P$, qui est irréductible
(puisqu'il est irréductible sur le corps éventuellement plus grand $E$)
et séparable (puisqu'il est séparable sur $E$).
Ainsi $k(a_0, \ldots, a_{d - 1}, \alpha)$ est séparable sur $k$
d'après le Lemme \ref{lemma-separably-generated-separable},
et nous concluons que $\alpha$ est séparable sur $k$, comme voulu.
\end{proof}

\begin{lemma}
\label{lemma-separable-elements}
Soit $E/k$ une extension de corps. Alors les éléments de $E$ séparables
sur $k$ forment une sous-extension de $E/k$.
\end{lemma}

\begin{proof}

Soient $\alpha, \beta \in E$ séparables sur $k$. Alors $\beta$ est séparable
sur $k(\alpha)$ d'après le Lemme \ref{lemma-separable-goes-up}.
Le Lemme \ref{lemma-separably-generated-separable} (appliqué avec $n = 2$,
$\alpha_1 = \alpha$ et $\alpha_2 = \beta$) montre que
$k(\alpha, \beta)$ est séparable sur $k$.

\end{proof}
\section{Indépendance linéaire des caractères}
\label{section-independence-characters}

\noindent
Voici l'énoncé.

\begin{lemma}
\label{lemma-independence-characters}
Soit $L$ un corps. Soit $G$ un monoïde, par exemple un groupe.
Soient
$\chi_1, \ldots, \chi_n : G \to L$ des homomorphismes de monoïdes deux à deux
distincts, où $L$ est considéré comme un monoïde pour la multiplication.
Alors $\chi_1, \ldots, \chi_n$ sont linéairement indépendants sur $L$ :
si $\lambda_1, \ldots, \lambda_n \in L$ ne sont pas tous nuls, alors
$\sum \lambda_i\chi_i(g) \not = 0$ pour un certain $g \in G$.
\end{lemma}

\begin{proof}
Si $n = 1$, cela résulte de $\chi_1(e) = 1$, où $e \in G$ est l'élément
neutre (identité). Nous démontrons le résultat par récurrence pour $n > 1$.
Supposons que $\lambda_1, \ldots, \lambda_n \in L$ ne soient pas tous nuls.
Si $\lambda_i = 0$ pour un certain $i$, le résultat découle de la récurrence sur $n$.
Puisque nous voulons montrer que $\sum \lambda_i\chi_i(g) \not = 0$
pour un certain $g \in G$, nous pouvons, après division par $-\lambda_n$,
supposer que $\lambda_n = -1$. La seule difficulté possible
surviendrait si
$$
\chi_n(g) = \sum\nolimits_{i = 1, \ldots, n - 1} \lambda_i\chi_i(g)
$$
pour tout $g \in G$. Fixons $h \in G$. Nous aurions alors également
\begin{align*}
\chi_n(h)\chi_n(g) & = \chi_n(hg) \\
& = \sum\nolimits_{i = 1, \ldots, n - 1} \lambda_i\chi_i(hg) \\
& = \sum\nolimits_{i = 1, \ldots, n - 1} \lambda_i\chi_i(h) \chi_i(g)
\end{align*}
En multipliant la relation précédente par $\chi_n(h)$ puis en soustrayant, on obtient
$$
0 = \sum\nolimits_{i = 1, \ldots, n - 1}
\lambda_i (\chi_n(h) - \chi_i(h)) \chi_i(g)
$$
pour tout $g \in G$. Comme $\lambda_i \not = 0$, la récurrence donne
$\chi_n(h) = \chi_i(h)$ pour tout $i$.
Le choix de $h$ étant arbitraire, nous en déduisons
que $\chi_i = \chi_n$ pour $i \leq n - 1$, ce qui contredit
l'hypothèse selon laquelle les caractères $\chi_i$ sont deux à deux distincts.
\end{proof}

\begin{lemma}
\label{lemma-sums-of-powers}
Soient $L$ un corps, $n \geq 1$ et $\alpha_1, \ldots, \alpha_n \in L$
des éléments deux à deux distincts de $L$. Alors il existe un
$e \geq 0$ tel que $\sum_{i = 1, \ldots, n} \alpha_i^e \not = 0$.
\end{lemma}

\begin{proof}
Appliquons l'indépendance linéaire des caractères
(Lemme \ref{lemma-independence-characters})
aux homomorphismes de monoïdes $\mathbf{Z}_{\geq 0} \to L$,
$e \mapsto \alpha_i^e$.
\end{proof}

\begin{lemma}
\label{lemma-independence-embeddings}
Soient $K/F$ et $L/F$ des extensions de corps. Soient
$\sigma_1, \ldots, \sigma_n : K \to L$ des morphismes de $F$-extensions
deux à deux distincts. Alors $\sigma_1, \ldots, \sigma_n$ sont linéairement
indépendants sur $L$ : si $\lambda_1, \ldots, \lambda_n \in L$ ne sont pas
tous nuls, alors $\sum \lambda_i\sigma_i(\alpha) \not = 0$
pour un certain $\alpha \in K$.
\end{lemma}

\begin{proof}
Appliquons le Lemme \ref{lemma-independence-characters} aux restrictions
des $\sigma_i$ aux groupes des unités.
\end{proof}

\begin{lemma}
\label{lemma-finite-separable-tensor-alg-closed}
Soient $K/F$ et $L/F$ des extensions de corps, avec
$K/F$ finie séparable et $L$ algébriquement clos.
Alors l'application
$$
K \otimes_F L
\longrightarrow
\prod\nolimits_{\sigma \in \Hom_F(K, L)} L,\quad
\alpha \otimes \beta \mapsto (\sigma(\alpha)\beta)_\sigma
$$
est un isomorphisme de $L$-algèbres.
\end{lemma}

\begin{proof}
Choisissons une base $\alpha_1, \ldots, \alpha_n$ de $K$ comme espace vectoriel sur $F$.
D'après le Lemme \ref{lemma-separable-equality} (et un petit argument omis), l'ensemble
$\Hom_F(K, L)$ possède $n$ éléments, disons $\sigma_1, \ldots, \sigma_n$.
En particulier, les deux membres ont la même dimension $n$ comme espaces
vectoriels sur $L$. Ainsi, si l'application n'est pas un isomorphisme, son
noyau est non nul. Autrement dit, il existerait des
$\mu_j \in L$, $j = 1, \ldots, n$, non tous nuls,
tels que $\sum \alpha_j \otimes \mu_j$ appartienne au noyau.
Autrement dit, $\sum \sigma_i(\alpha_j)\mu_j = 0$ pour tout $i$.
Cela signifierait que la matrice $n \times n$ de coefficients
$\sigma_i(\alpha_j)$ n'est pas inversible. Nous pouvons donc trouver
$\lambda_1, \ldots, \lambda_n \in L$, non tous nuls,
tels que $\sum \lambda_i\sigma_i(\alpha_j) = 0$ pour tout $j$.
Or tout élément $\alpha \in K$ peut s'écrire
$\alpha = \sum \beta_j \alpha_j$ avec $\beta_j \in F$, et nous obtiendrions
$$
\sum \lambda_i\sigma_i(\alpha) =
\sum \lambda_i\sigma_i(\sum \beta_j \alpha_j) =
\sum \beta_j \sum \lambda_i\sigma_i(\alpha_j) = 0
$$
ce qui contredit le Lemme \ref{lemma-independence-embeddings}.
\end{proof}







\section{Extensions purement inséparables}
\label{section-purely-inseparable}

\noindent
Les extensions purement inséparables sont à l'opposé des extensions
séparables définies à la section précédente. Elles n'apparaissent
qu'en caractéristique positive.

\begin{definition}
\label{definition-purely-inseparable}
Soit $F$ un corps de caractéristique $p > 0$. Soit $K/F$ une extension.
\begin{enumerate}
\item Un élément $\alpha \in K$ est dit {\it purement inséparable} sur $F$
s'il existe une puissance $q$ de $p$ telle que $\alpha^q \in F$.
\item L'extension $K/F$ est dite {\it purement inséparable}
si et seulement si tout élément de $K$ est purement inséparable sur $F$.
\end{enumerate}
Si $L/M$ est une extension de corps (sans condition sur la caractéristique
de $M$), nous dirons que l'extension est {\it purement inséparable}
si $L = M$, ou si la caractéristique de $M$ est un nombre premier $p$
et que $L/M$ est purement inséparable au sens défini ci-dessus.
\end{definition}

\noindent
Observons qu'une extension purement inséparable est nécessairement algébrique.
Soit $F$ un corps de caractéristique $p > 0$.
On obtient un exemple d'extension purement inséparable en adjoignant
la racine $p$-ième d'un élément $t \in F$ qui n'en possède pas encore. En effet,
le lemme ci-dessous montre que $P = x^p - t$ est irréductible, et donc
$$
K = F[x]/(P) = F[t^{1/p}]
$$
est un corps. De plus, $K$ est purement inséparable sur $F$, car tout élément
$$
a_0 + a_1t^{1/p} + \ldots + a_{p - 1}t^{(p - 1)/p},\quad a_i \in F
$$
de $K$ a pour puissance $p$-ième
$$
(a_0 + a_1t^{1/p} + \ldots + a_{p - 1}t^{(p - 1)/p})^p =
a_0^p + a_1^p t + \ldots + a_{p - 1}^pt^{p - 1} \in F
$$
Cette situation se présente pour le corps
$\mathbf{F}_p(t)$ des fonctions rationnelles sur $\mathbf{F}_p$.

\begin{lemma}
\label{lemma-take-pth-root}
Soit $p$ un nombre premier. Soit $F$ un corps de caractéristique $p$.
Soit $t \in F$ un élément qui n'a pas de racine $p$-ième dans $F$.
Alors le polynôme $x^p - t$ est irréductible sur $F$.
\end{lemma}

\begin{proof}
Pour le voir, supposons donnée une factorisation
$x^p - t = f g$. En dérivant, on obtient $f' g + f g' = 0$.
Notons que ni $f' = 0$ ni $g' = 0$, puisque les degrés de $f$ et de $g$
sont strictement inférieurs à $p$. En outre, $\deg(f') < \deg(f)$ et $\deg(g') < \deg(g)$.
Nous en déduisons que $f$ et $g$ ont un facteur commun. Ainsi, si $x^p - t$
est réductible, il est de la forme $x^p - t = c f^n$ pour un certain polynôme
irréductible $f$, un certain $c \in F^*$ et un certain $n > 1$. Comme $p$ est premier,
cela entraîne $n = p$ et $f$ linéaire, ce qui impliquerait que $x^p - t$ possède
une racine dans $F$. Contradiction.
\end{proof}

\noindent
Nous verrons que l'extraction de racines $p$-ièmes est une opération très importante
en caractéristique $p$.

\begin{lemma}
\label{lemma-purely-inseparable-permanence}
Soient $E/k$ et $F/E$ des extensions de corps purement inséparables. Alors $F/k$
est une extension de corps purement inséparable.
\end{lemma}

\begin{proof}
Si la caractéristique est nulle, les deux extensions sont triviales et il n'y a
rien à démontrer. Supposons donc que la caractéristique de $k$ soit un nombre
premier $p$. Choisissons $\alpha \in F$. Alors
$\alpha^q \in E$ pour une certaine puissance de $p$, notée $q$. Dès lors, $(\alpha^q)^{q'} \in k$
pour une certaine puissance de $p$, notée $q'$. Par conséquent, $\alpha^{qq'} \in k$.
\end{proof}

\begin{lemma}
\label{lemma-purely-inseparable-elements}
Soit $E/k$ une extension de corps. Alors les éléments de $E$ purement inséparables
sur $k$ forment une sous-extension de $E/k$.
\end{lemma}

\begin{proof}
Soit $p$ la caractéristique de $k$.
Si celle-ci est nulle, l'assertion est immédiate, car tout élément purement
inséparable appartient alors au corps de base. Supposons donc cette caractéristique
positive.
Soient $\alpha, \beta \in E$ purement inséparables sur $k$. Écrivons
$\alpha^q \in k$ et $\beta^{q'} \in k$ pour certaines puissances de $p$, notées $q, q'$.
Si $q''$ est une puissance de $p$, alors
$(\alpha + \beta)^{q''} = \alpha^{q''} + \beta^{q''}$.
Par conséquent, si $q'' \geq q, q'$, nous en déduisons que $\alpha + \beta$
est purement inséparable sur $k$. Il en va de même de la différence,
du produit et du quotient de $\alpha$ et $\beta$.
\end{proof}

\begin{lemma}
\label{lemma-finite-purely-inseparable}
Soit $E/F$ une extension de corps finie purement inséparable de
caractéristique $p > 0$. Alors il existe une suite d'éléments
$\alpha_1, \ldots, \alpha_n \in E$ qui donne une tour
de corps
$$
E = F(\alpha_1, \ldots, \alpha_n) \supset
F(\alpha_1, \ldots, \alpha_{n - 1}) \supset
\ldots
\supset F(\alpha_1) \supset F
$$
telle que chaque extension intermédiaire soit de degré $p$ et s'obtienne
par adjonction d'une racine $p$-ième. Plus précisément,
$\alpha_i^p \in F(\alpha_1, \ldots, \alpha_{i - 1})$
est un élément qui ne possède pas de racine $p$-ième dans
$F(\alpha_1, \ldots, \alpha_{i - 1})$, pour $i = 1, \ldots, n$.
\end{lemma}

\begin{proof}
Raisonnons par récurrence sur le degré de $E/F$. Si ce degré vaut $1$,
le résultat est clair (avec $n = 0$). Sinon, choisissons
$\alpha \in E$, $\alpha \not \in F$. Écrivons $\alpha^{p^r} \in F$ pour un certain
$r > 0$. Prenons $r$ minimal et remplaçons $\alpha$ par $\alpha^{p^{r - 1}}$.
Alors $\alpha \not \in F$, mais $\alpha^p \in F$. Ainsi $t = \alpha^p$ n'est pas
une puissance $p$-ième dans $F$ (car cela entraînerait $\alpha \in F$, voir le
Lemme \ref{lemma-nr-roots-unchanged} ou sa démonstration).
Ainsi, $F \subset F(\alpha)$ est une sous-extension de degré $p$
(Lemme \ref{lemma-take-pth-root}). Par récurrence, nous trouvons
$\alpha_1, \ldots, \alpha_n \in E$ qui engendrent $E/F(\alpha)$
et satisfont aux conclusions du lemme.
La suite $\alpha, \alpha_1, \ldots, \alpha_n$ convient alors
pour l'extension $E/F$.
\end{proof}

\begin{lemma}
\label{lemma-separable-first}
\begin{slogan}
Toute extension algébrique de corps est, de manière unique, une extension
séparable suivie d'une extension purement inséparable.
\end{slogan}
Soit $E/F$ une extension algébrique de corps. Il existe une unique sous-extension
$E/E_{sep}/F$ telle que $E_{sep}/F$ soit séparable et que $E/E_{sep}$ soit
purement inséparable.
\end{lemma}

\begin{proof}
Si la caractéristique est nulle, posons $E_{sep} = E$. Supposons désormais la caractéristique
égale à $p > 0$. Soit $E_{sep}$ l'ensemble des éléments de $E$ qui sont séparables
sur $F$. C'est une sous-extension d'après le Lemme \ref{lemma-separable-elements},
et bien entendu $E_{sep}$ est séparable sur $F$. Étant donné $\alpha$ dans $E$,
il existe une puissance de $p$, notée $q$, telle que $\alpha^q$ soit séparable sur $F$.
En effet, $q$ est la puissance de $p$ telle que le polynôme minimal de
$\alpha$ soit de la forme $P(x^q)$, où $P$ est un polynôme séparable, voir le
Lemme \ref{lemma-irreducible-polynomials}. Ainsi $E/E_{sep}$ est purement
inséparable. L'unicité est claire.
\end{proof}

\begin{definition}
\label{definition-insep-degree}
Soit $E/F$ une extension algébrique de corps. Soit $E_{sep}$ la sous-extension
obtenue au Lemme \ref{lemma-separable-first}.
\begin{enumerate}
\item Le degré $[E_{sep} : F]$ est appelé le {\it degré séparable}
de l'extension. Notation : $[E : F]_s$.
\item Le degré $[E : E_{sep}]$ est appelé le {\it degré inséparable},
ou le {\it degré d'inséparabilité} de l'extension.
Notation : $[E : F]_i$.
\end{enumerate}
\end{definition}

\noindent
Bien entendu, en caractéristique $0$, on a $[E : F] = [E : F]_s$ et
$[E : F]_i = 1$. Par multiplicativité
(Lemme \ref{lemma-multiplicativity-degrees}), on a
$$
[E : F] = [E : F]_s [E : F]_i
$$
même lorsque certains de ces degrés sont infinis. En fait, le degré séparable
et le degré inséparable sont eux aussi multiplicatifs (voir le
Lemme \ref{lemma-multiplicativity-all-degrees}).

\begin{lemma}
\label{lemma-separable-degree}
Soit $K/F$ une extension finie. Soit $\overline{F}$ une clôture algébrique
de $F$. Alors $[K : F]_s = |\Mor_F(K, \overline{F})|$.
\end{lemma}

\begin{proof}
Démontrons d'abord ce résultat lorsque $K/F$ est purement inséparable. Nous affirmons
que, dans ce cas, il existe un unique $F$-plongement $K \to \overline{F}$. On peut le
voir en choisissant une suite d'éléments $\alpha_1, \ldots, \alpha_n \in K$
comme au Lemme \ref{lemma-finite-purely-inseparable}. Le polynôme irréductible
de $\alpha_i$ sur $F(\alpha_1, \ldots, \alpha_{i - 1})$ est $x^p - \alpha_i^p$.
En appliquant le Lemme \ref{lemma-count-embeddings-explicitly}, on voit que
$|\Mor_F(K, \overline{F})| = 1$. D'autre part, $[K : F]_s = 1$
dans ce cas, d'où l'égalité.

\medskip\noindent
Revenons à une extension finie générale $K/F$. Dans ce cas,
choisissons $F \subset K_s \subset K$ comme au Lemme \ref{lemma-separable-first}.
D'après le Lemme \ref{lemma-separable-equality}, on a
$|\Mor_F(K_s, \overline{F})| = [K_s : F] = [K : F]_s$.
D'autre part, tout morphisme de corps $\sigma' : K_s \to \overline{F}$
se prolonge de manière unique en un morphisme de corps $\sigma : K \to \overline{F}$, d'après le
résultat du paragraphe précédent. Autrement dit,
$|\Mor_F(K, \overline{F})| = |\Mor_F(K_s, \overline{F})|$,
ce qui achève la démonstration.
\end{proof}

\begin{lemma}[Multiplicativité]
\label{lemma-multiplicativity-all-degrees}
Supposons donnée une tour d'extensions algébriques de corps $K/E/F$. Alors
$$
[K : F]_s = [K : E]_s [E : F]_s
\quad\text{et}\quad
[K : F]_i = [K : E]_i [E : F]_i
$$
\end{lemma}

\begin{proof}
Démontrons d'abord le résultat lorsque $K$ est fini sur $F$. Comme la
multiplicativité vaut pour le degré usuel (d'après le
Lemme \ref{lemma-multiplicativity-degrees}), il suffit de démontrer
l'une des deux formules. D'après le Lemme \ref{lemma-separable-degree}, on a
$[K : F]_s = |\Mor_F(K, \overline{F})|$. Par le même lemme,
pour tout $\sigma \in \Mor_F(E, \overline{F})$, le nombre de prolongements
de $\sigma$ en une application $\tau : K \to \overline{F}$ est $[K : E]_s$.
En effet, via $E \cong \sigma(E) \subset \overline{F}$, nous pouvons considérer
$\overline{F}$ comme une clôture algébrique de $E$. En combinant ceci avec le
fait qu'il existe $[E : F]_s = |\Mor_F(E, \overline{F})|$ choix
pour $\sigma$, on obtient le résultat.

\medskip\noindent
Nous omettons la démonstration lorsque les extensions sont infinies.
\end{proof}









\section{Extensions normales}
\label{section-normal}

\noindent
Soit $P \in F[x]$ un polynôme non constant sur un corps $F$. Nous dirons que $P$
{\it se décompose complètement en facteurs linéaires sur $F$} ou
{\it est complètement décomposé sur $F$} s'il existe
$c \in F^*$, $n \geq 1$, $\alpha_1, \ldots, \alpha_n \in F$ tels que
$$
P = c(x - \alpha_1) \ldots (x - \alpha_n)
$$
dans $F[x]$. Les extensions normales sont définies comme suit.

\begin{definition}
\label{definition-normal}
Soit $E/F$ une extension algébrique de corps. Nous dirons que $E$ est {\it normale}
sur $F$ si, pour tout $\alpha \in E$, le polynôme minimal $P$
de $\alpha$ sur $F$ se décompose complètement en facteurs linéaires sur $E$.
\end{definition}

\noindent
Comme pour les extensions séparables, l'établissement des propriétés
fondamentales de cette notion demande un peu de travail.

\begin{lemma}
\label{lemma-normal-goes-up}
Soit $K/E/F$ une tour d'extensions algébriques de corps.
Si $K$ est normale sur $F$, alors $K$ est normale sur $E$.
\end{lemma}

\begin{proof}
Soit $\alpha \in K$. Soit $P$ le polynôme minimal de $\alpha$ sur $F$.
Soit $Q$ le polynôme minimal de $\alpha$ sur $E$.
Alors $Q$ divise $P$ dans l'anneau de polynômes $E[x]$, disons $P = QR$.
Par conséquent, si $P$ est complètement décomposé sur $K$, il en va de même de $Q$.
\end{proof}

\begin{lemma}
\label{lemma-intersect-normal}
Soit $F$ un corps. Soit $M/F$ une extension algébrique. Soient
$M/E_i/F$, $i \in I$, une famille non vide de sous-extensions telles que
$E_i/F$ soit normale. Alors $\bigcap E_i$ est normale sur $F$.
\end{lemma}

\begin{proof}
Cela résulte directement des définitions.
\end{proof}

\begin{lemma}
\label{lemma-separable-first-normal}
Soit $E/F$ une extension algébrique de corps normale. Alors la sous-extension
$E/E_{sep}/F$ du Lemme \ref{lemma-separable-first} est normale.
\end{lemma}

\begin{proof}
Si la caractéristique est nulle, alors $E_{sep} = E$, et le résultat
est clair. Si la caractéristique vaut $p > 0$, alors $E_{sep}$
est l'ensemble des éléments de $E$ qui sont séparables sur $F$.
Alors, si $\alpha \in E_{sep}$ a pour polynôme minimal $P$,
écrivons $P = c(x - \alpha)(x - \alpha_2) \ldots (x - \alpha_d)$
avec $\alpha_2, \ldots, \alpha_d \in E$. Puisque
$P$ est un polynôme séparable et que $\alpha_i$
est une racine de $P$, nous en déduisons $\alpha_i \in E_{sep}$, comme voulu.
\end{proof}

\begin{lemma}
\label{lemma-characterize-normal}
Soit $E/F$ une extension algébrique de corps. Soit $\overline{F}$ une
clôture algébrique de $F$. Les assertions suivantes sont équivalentes :
\begin{enumerate}
\item $E$ est normale sur $F$ ;
\item pour toute paire $\sigma, \sigma' \in \Mor_F(E, \overline{F})$, on
a $\sigma(E) = \sigma'(E)$.
\end{enumerate}
\end{lemma}

\begin{proof}
Soit $\mathcal{P}$ l'ensemble de tous les polynômes minimaux sur $F$
de tous les éléments de $E$. Posons
$$
T =
\{\beta \in \overline{F} \mid P(\beta) = 0\text{ pour un certain }P \in \mathcal{P}\}
$$
Il est clair que, si $E$ est normale sur $F$, alors $\sigma(E) = T$
pour tout $\sigma \in \Mor_F(E, \overline{F})$. Nous voyons donc que (1)
implique (2).

\medskip\noindent
Réciproquement, supposons (2). Choisissons $\beta \in T$.
Nous pouvons trouver $\alpha \in E$ dont le polynôme minimal
$P \in \mathcal{P}$ annule $\beta$. Puisque $F(\alpha) = F[x]/(P)$,
nous pouvons trouver un élément $\sigma_0 \in \Mor_F(F(\alpha), \overline{F})$ qui envoie
$\alpha$ sur $\beta$. D'après le Lemme \ref{lemma-map-into-algebraic-closure},
nous pouvons prolonger $\sigma_0$ en un morphisme $\sigma \in \Mor_F(E, \overline{F})$.
Nous voyons ainsi que $\beta$ appartient à l'image commune de tous les plongements
$\sigma : E \to \overline{F}$. Il s'ensuit que $\sigma(E) = T$
pour tout $\sigma$. Fixons un $\sigma$. Soit maintenant $P \in \mathcal{P}$. Alors nous
pouvons écrire
$$
P = (x - \beta_1) \ldots (x - \beta_n)
$$
pour un certain $n$ et des $\beta_i \in \overline{F}$, d'après le
Lemme \ref{lemma-algebraically-closed}. Observons que $\beta_i \in T$.
Ainsi $\beta_i = \sigma(\alpha_i)$ pour un certain $\alpha_i \in E$. Par conséquent,
$P = (x - \alpha_1) \ldots (x - \alpha_n)$ se décompose complètement sur $E$.
Cela achève la démonstration.
\end{proof}

\begin{lemma}
\label{lemma-normally-generated}
Soit $E/F$ une extension algébrique de corps.
Si $E$ est engendrée par les éléments $\alpha_i \in E$, $i \in I$,
sur $F$, et si, pour tout $i$, le polynôme minimal
de $\alpha_i$ sur $F$ se décompose complètement dans $E$, alors
$E/F$ est normale.
\end{lemma}

\begin{proof}
Soit $P_i$ le polynôme minimal de $\alpha_i$ sur $F$.
Soient $\alpha_i = \alpha_{i, 1}, \alpha_{i, 2}, \ldots, \alpha_{i, d_i}$
les racines de $P_i$ sur $E$. Étant donnés deux plongements
$\sigma, \sigma' : E \to \overline{F}$ au-dessus de $F$, on voit que
$$
\{\sigma(\alpha_{i, 1}), \ldots, \sigma(\alpha_{i, d_i})\} =
\{\sigma'(\alpha_{i, 1}), \ldots, \sigma'(\alpha_{i, d_i})\}
$$
car les deux membres sont égaux à l'ensemble des racines de $P_i$
dans $\overline{F}$. Les éléments $\alpha_{i, j}$
engendrent $E$ sur $F$, et l'on obtient $\sigma(E) = \sigma'(E)$.
Ainsi $E/F$ est normale d'après le Lemme \ref{lemma-characterize-normal}.
\end{proof}

\begin{lemma}
\label{lemma-lift-maps}
Soit $L/M/K$ une tour d'extensions algébriques.
\begin{enumerate}
\item Si $M/K$ est normale, alors tout automorphisme $\tau$ de $L/K$
induit un automorphisme $\tau|_M : M \to M$.
\item Si $L/K$ est normale, alors tout morphisme de $K$-algèbres $\sigma : M \to L$
se prolonge en un automorphisme de $L$.
\end{enumerate}
\end{lemma}

\begin{proof}
Choisissons une clôture algébrique $\overline{L}$ de $L$
(Théorème \ref{theorem-existence-algebraic-closure}).

\medskip\noindent
Soit $\tau$ comme en (1). Alors $\tau(M) = M$ comme sous-corps de $\overline{L}$
d'après le Lemme \ref{lemma-characterize-normal}, et donc
$\tau|_M : M \to M$ est un automorphisme.

\medskip\noindent
Soit $\sigma : M \to L$ comme en (2).
D'après le Lemme \ref{lemma-map-into-algebraic-closure},
nous pouvons prolonger $\sigma$ en une application
$\tau : L \to \overline{L}$, c'est-à-dire de telle sorte que le diagramme
$$
\xymatrix{
L \ar[r]_\tau & \overline{L} \\
M \ar[u] \ar[ru]_\sigma & K \ar[l] \ar[u]
}
$$
soit commutatif. D'après le Lemme \ref{lemma-characterize-normal}, on voit que
$\tau(L) = L$. Ainsi $\tau : L \to L$ est un automorphisme qui
prolonge $\sigma$.
\end{proof}

\begin{definition}
\label{definition-automorphisms}
Soit $E/F$ une extension de corps. Alors $\text{Aut}(E/F)$ ou
$\text{Aut}_F(E)$ désigne le groupe des automorphismes de $E$ comme objet
de la catégorie des $F$-extensions. Les éléments de $\text{Aut}(E/F)$
sont appelés {\it automorphismes de $E$ sur $F$} ou
{\it automorphismes de $E/F$}.
\end{definition}

\noindent
Voici une caractérisation des extensions normales en termes d'automorphismes.

\begin{lemma}
\label{lemma-normal-and-automorphisms}
Soit $E/F$ une extension finie. On a
$$
|\text{Aut}(E/F)| \leq [E : F]_s
$$
avec égalité si et seulement si $E$ est normale sur $F$.
\end{lemma}

\begin{proof}
Choisissons une clôture algébrique $\overline{F}$ de $F$. Rappelons que
$[E : F]_s = |\Mor_F(E, \overline{F})|$. Choisissons un élément
$\sigma_0 \in \Mor_F(E, \overline{F})$. Alors l'application
$$
\text{Aut}(E/F) \longrightarrow \Mor_F(E, \overline{F}),\quad
\tau \longmapsto \sigma_0 \circ \tau
$$
est injective. D'où l'inégalité. S'il y a égalité, tout
$\sigma \in \Mor_F(E, \overline{F})$ s'obtient en précomposant
$\sigma_0$ par un automorphisme. Ainsi $\sigma(E) = \sigma_0(E)$.
Donc $E$ est normale sur $F$ d'après le Lemme \ref{lemma-characterize-normal}.

\medskip\noindent
Réciproquement, supposons que $E/F$ est normale. Alors, d'après le
Lemme \ref{lemma-characterize-normal}, on a $\sigma(E) = \sigma_0(E)$
pour tout $\sigma \in \Mor_F(E, \overline{F})$.
On obtient donc un automorphisme de $E$ sur $F$ en posant
$\tau = \sigma_0^{-1} \circ \sigma$. Par conséquent, l'application ci-dessus
est surjective.
\end{proof}

\begin{lemma}
\label{lemma-normal-embeddings-differ-by-aut}
Soit $L/K$ une extension algébrique normale de corps.
Soit $E/K$ une extension de corps. Alors, soit
il n'existe aucun $K$-plongement de $L$ dans $E$, soit
il en existe un $\tau : L \to E$ et tout autre
est de la forme $\tau \circ \sigma$, où $\sigma \in \text{Aut}(L/K)$.
\end{lemma}

\begin{proof}
Étant donné $\tau$, remplaçons $L$ par $\tau(L) \subset E$ et appliquons le
Lemme \ref{lemma-lift-maps}.
\end{proof}





\section{Corps de décomposition}
\label{section-splitting-fieds}

\noindent
Le lemme suivant est un outil utile pour construire des extensions normales de corps.

\begin{lemma}
\label{lemma-splitting-field}
Soit $F$ un corps. Soit $P \in F[x]$ un polynôme non constant.
Il existe une plus petite extension de corps $E/F$ telle que $P$
soit scindé sur $E$. De plus, l'extension de corps $E/F$ est normale
et unique à isomorphisme près, cet isomorphisme n'étant pas nécessairement unique.
\end{lemma}

\begin{proof}
Choisissons une clôture algébrique $\overline{F}$. Nous pouvons alors écrire
$P = c (x - \beta_1) \ldots (x - \beta_n)$ dans $\overline{F}[x]$, voir le
Lemme \ref{lemma-algebraically-closed}. Notons que $c \in F^*$. Posons
$E = F(\beta_1, \ldots, \beta_n)$. Il est alors clair que $E$ est
minimale sous la condition que $P$ soit scindé sur $E$.

\medskip\noindent
Soit ensuite $E'$ une autre extension minimale de $F$ telle que
$P$ soit scindé sur $E'$. Écrivons
$P = c (x - \alpha_1) \ldots (x - \alpha_n)$ avec $c \in F$ et
$\alpha_i \in E'$. La minimalité entraîne de nouveau que
$E' = F(\alpha_1, \ldots, \alpha_n)$. De plus, si nous choisissons
un $\sigma : E' \to \overline{F}$ quelconque
(Lemme \ref{lemma-map-into-algebraic-closure}),
nous voyons immédiatement que $\sigma(\alpha_i) = \beta_{\tau(i)}$
pour une certaine permutation $\tau : \{1, \ldots, n\} \to \{1, \ldots, n\}$.
Ainsi $\sigma(E') = E$. Cela implique que $E'$ est une extension normale
de $F$ d'après le Lemme \ref{lemma-characterize-normal},
et que $E \cong E'$ comme extensions de $F$, ce qui achève la démonstration.
\end{proof}

\begin{definition}
\label{definition-splitting-field}
Soit $F$ un corps. Soit $P \in F[x]$ un polynôme non constant.
L'extension de corps $E/F$ construite au Lemme \ref{lemma-splitting-field}
est appelée le {\it corps de décomposition de $P$ sur $F$}.
\end{definition}

\begin{lemma}
\label{lemma-normal-closure}
\begin{slogan}
Existence de la clôture normale des extensions finies de corps.
\end{slogan}
Soit $E/F$ une extension finie de corps. Il existe une unique
plus petite extension finie $K/E$ telle que $K$ soit normale sur $F$.
\end{lemma}

\begin{proof}
Choisissons des générateurs $\alpha_1, \ldots, \alpha_n$ de $E$ sur $F$.
Soient $P_1, \ldots, P_n$ les polynômes minimaux de
$\alpha_1, \ldots, \alpha_n$ sur $F$. Posons $P = P_1 \ldots P_n$.
Observons que $(x - \alpha_1) \ldots (x - \alpha_n)$ divise $P$, puisque
chaque $(x - \alpha_i)$ divise $P_i$. Écrivons
$P = (x - \alpha_1) \ldots (x - \alpha_n)Q$.
Soit $K/E$ le corps de décomposition de $P$ sur $E$.
Nous affirmons que $K$ est également le corps de décomposition de $P$ sur $F$
(ce qui implique que $K$ est normale sur $F$).
Cela est clair, car $K/E$ est engendrée par les racines de
$Q$ sur $E$, tandis que $E$ est engendrée par les racines de
$(x - \alpha_1) \ldots (x - \alpha_n)$ sur $F$ ; ainsi,
$K$ est engendrée par les racines de $P$ sur $F$.

\medskip\noindent
Unicité. Supposons que $K'/E$ soit une seconde extension minimale telle que
$K'/F$ soit normale. Choisissons une clôture algébrique $\overline{F}$ et un
plongement $\sigma_0 : E \to \overline{F}$. D'après le
Lemme \ref{lemma-map-into-algebraic-closure},
nous pouvons prolonger $\sigma_0$ en $\sigma : K \to \overline{F}$ et en
$\sigma' : K' \to \overline{F}$.
D'après le Lemme \ref{lemma-intersect-normal}, on voit que
$\sigma(K) \cap \sigma'(K')$ est normale sur $F$.
La minimalité donne $\sigma(K) = \sigma'(K')$.
Ainsi $\sigma^{-1} \circ \sigma' : K' \to K$ fournit un isomorphisme
d'extensions de $E$.
\end{proof}

\begin{definition}
\label{definition-normal-closure}
Soit $E/F$ une extension finie de corps. L'extension de corps $K/E$
construite au Lemme \ref{lemma-normal-closure}
est appelée la {\it clôture normale de $E$ sur $F$}.
\end{definition}

\noindent
On peut construire la clôture normale à l'intérieur de toute extension normale donnée.

\begin{lemma}
\label{lemma-normal-closure-inside-normal}
Soit $L/K$ une extension algébrique normale.
\begin{enumerate}
\item Si $L/M/K$ est une sous-extension avec $M/K$ finie, alors il existe
une tour $L/M'/M/K$ telle que $M'/K$ soit finie et normale.
\item Si $L/M'/M/K$ est une tour telle que $M/K$ soit normale et $M'/M$ finie,
alors il existe une tour $L/M''/M'/M/K$ telle que $M''/M$ soit
finie et $M''/K$ normale.
\end{enumerate}
\end{lemma}

\begin{proof}
Démonstration de (1). Soit $M'$ la plus petite sous-extension de $L/K$ contenant $M$
qui soit normale sur $K$. D'après le Lemme \ref{lemma-normal-closure},
c'est la clôture normale de $M/K$, et elle est finie sur $K$.

\medskip\noindent
Démonstration de (2). Soient $\alpha_1, \ldots, \alpha_n \in M'$ des générateurs de $M'$ sur $M$.
Soient $P_1, \ldots, P_n$ les polynômes minimaux de
$\alpha_1, \ldots, \alpha_n$ sur $K$. Soient $\alpha_{i, j}$ les racines
de $P_i$ dans $L$. Posons $M''  = M(\alpha_{i, j})$. Il résulte du
Lemme \ref{lemma-normally-generated}
(appliqué à l'ensemble de générateurs $M \cup \{\alpha_{i, j}\}$)
que $M''$ est normale sur $K$.
\end{proof}

\noindent
Le lemme suivant peut parfois servir à démontrer des propriétés
de la clôture normale.

\begin{lemma}
\label{lemma-normal-closure-tensor-product}
Soit $L/K$ une extension finie. Soit $M/L$ la clôture
normale de $L$ sur $K$. Alors il existe une application surjective
$$
L \otimes_K L \otimes_K \ldots \otimes_K L \longrightarrow M
$$
de $K$-algèbres, où le nombre de facteurs tensoriels peut être pris égal à
$[L : K]_s \leq [L : K]$.
\end{lemma}

\begin{proof}
Choisissons une clôture algébrique $\overline{K}$ de $K$.
Posons $n = [L : K]_s = |\Mor_K(L, \overline{K})|$, l'égalité résultant du
Lemme \ref{lemma-separable-degree}.
Écrivons $\Mor_K(L, \overline{K}) = \{\sigma_1, \ldots, \sigma_n\}$.
Soit $M' \subset \overline{K}$ la $K$-sous-algèbre
engendrée par $\sigma_i(L)$, $i = 1, \ldots, n$. Alors $M'$ est
un corps, car toute $K$-sous-algèbre de $\overline{K}$ est un corps.
Tout morphisme de $K$-algèbres de $M'$ vers $\overline{K}$ permute les $\sigma_i$,
et envoie donc le sous-corps $M'$ sur $M'$. Par construction, le corps $M'$
est engendré par les conjugués d'éléments de $\sigma_1(L)$. Cela étant,
il résulte du Lemme \ref{lemma-characterize-normal}
que $M'$ est normale sur $K$ et qu'elle est la
plus petite sous-extension normale de $\overline{K}$ contenant
$\sigma_1(L)$. Par unicité de la clôture normale, on a $M \cong M'$.
Enfin, il existe une application surjective
$$
L \otimes_K L \otimes_K \ldots \otimes_K L \longrightarrow M',
\quad
\lambda_1 \otimes \ldots \otimes \lambda_n \longmapsto
\sigma_1(\lambda_1) \ldots \sigma_n(\lambda_n)
$$
et notons que $n \leq [L : K]$ par définition.
\end{proof}












\section{Racines de l’unité}
\label{section-roots-of-1}

\noindent
Soit $F$ un corps. Pour un entier $n \geq 1$, posons
$$
\mu_n(F) = \{\zeta \in F \mid \zeta^n = 1\}
$$
On appelle cet ensemble le {\it groupe des racines $n$-ièmes de l'unité} ; ses
éléments sont les {\it racines $n$-ièmes de $1$}. C'est un groupe abélien pour la multiplication,
dont l'élément neutre est $1$.
Observons que, dans un corps, le nombre de racines d'un polynôme de degré $d$
est toujours au plus $d$. Ainsi $|\mu_n(F)| \leq n$,
puisque cet ensemble est défini par une équation polynomiale de degré $n$.
Bien entendu, tout élément de $\mu_n(F)$ est d'ordre divisant $n$.
En outre, les sous-groupes
$$
\mu_d(F) \subset \mu_n(F),\quad d | n
$$
ont chacun au plus $d$ éléments. Cela implique que $\mu_n(F)$ est cyclique.

\begin{lemma}
\label{lemma-cyclic}
Soit $A$ un groupe abélien d'exposant divisant $n$ tel que
$\{x \in A \mid dx = 0\}$ soit de cardinal au plus $d$ pour tout $d | n$.
Alors $A$ est cyclique d'ordre divisant $n$.
\end{lemma}

\begin{proof}
Si $A$ est trivial, il n'y a rien à démontrer. Supposons désormais que $A$
ne soit pas trivial.
Les conditions impliquent que $|A| \leq n$ ; en particulier, $A$ est fini.
La structure des groupes abéliens finis montre que
$A = \mathbf{Z}/e_1\mathbf{Z} \oplus \ldots \oplus \mathbf{Z}/e_r\mathbf{Z}$
pour des entiers $1 < e_1 | e_2 | \ldots | e_r$. Cela impliquerait
que $\{x \in A \mid e_1 x = 0\}$ est de cardinal $e_1^r$. Par conséquent,
$r = 1$.
\end{proof}

\noindent
Appliqué au corps $\mathbf{F}_p$, ceci donne le célèbre résultat
selon lequel le groupe $(\mathbf{Z}/p\mathbf{Z})^*$ est cyclique. Nous reviendrons sur
cette question dans la section consacrée aux corps finis.

\medskip\noindent
Une autre observation est souvent utile : si $F$ est de caractéristique
$p > 0$, alors $\mu_{p^n}(F) = \{1\}$. Cela tient à ce que l'élévation
à la puissance $p$ est une application injective sur les corps de caractéristique $p$,
comme nous l'avons vu dans la démonstration du Lemme \ref{lemma-nr-roots-unchanged}.
(Bien entendu, cela résulte aussi de l'énoncé même de ce lemme.)






\section{Corps finis}
\label{section-finite}

\noindent
Soit $F$ un corps fini. Il est clair que $F$ est de caractéristique positive,
car il ne peut exister d'injection $\mathbf{Q} \to F$. Supposons que la caractéristique
de $F$ soit $p$. L'extension $\mathbf{F}_p \subset F$ est finie.
Nous voyons ainsi que $F$ possède $q = p^f$ éléments pour un certain $f \geq 1$.

\medskip\noindent
Considérons le groupe des unités $F^*$. C'est un groupe abélien
fini, il possède donc un certain exposant $e$. Alors $F^* = \mu_e(F)$, et la
discussion de la Section \ref{section-roots-of-1} montre que $F^*$
est un groupe cyclique d'ordre $q - 1$. (A posteriori, il en résulte aussi que
$e = q - 1$.) En particulier, si $\alpha \in F^*$ est un générateur,
alors il est clair que
$$
F = \mathbf{F}_p(\alpha)
$$
Autrement dit, l'extension $F/\mathbf{F}_p$ est engendrée par un seul
élément. Bien entendu, il en va de même de toute extension de corps finis
$E/F$ (car $E$ est déjà engendré par un seul élément sur
le corps premier).





\section{Éléments primitifs}
\label{section-primitive-element}

\noindent
Soit $E/F$ une extension finie de corps. Un élément $\alpha \in E$
est appelé {\it élément primitif de $E$ sur $F$} si $E = F(\alpha)$.

\begin{lemma}[Élément primitif]
\label{lemma-primitive-element}
Soit $E/F$ une extension finie de corps. Les assertions suivantes sont équivalentes :
\begin{enumerate}
\item il existe un élément primitif de $E$ sur $F$ ;
\item il existe un nombre fini de sous-extensions $E/K/F$.
\end{enumerate}
De plus, (1) et (2) sont vérifiées si $E/F$ est séparable.
\end{lemma}

\begin{proof}
Soit $\alpha \in E$ un élément primitif. Soit $P$ le polynôme minimal
de $\alpha$ sur $F$. Soit ensuite $E/K/F$ une sous-extension.
Soit $Q$ le polynôme minimal de $\alpha$ sur $K$. Observons que
$\deg(Q) = [E : K]$. En écrivant $Q = x^d + \sum_{i < d} a_i x^i$, nous affirmons que
$K$ est égal à $L = F(a_0, \ldots, a_{d - 1})$. En effet, $\alpha$ est de degré
$d$ sur $L$ et $L \subset K$. Ainsi $[E : L] = [E : K]$, d'où
$[K : L] = 1$, c'est-à-dire $K = L$.
Il suffit donc de montrer qu'il n'existe qu'un nombre fini de possibilités
pour le polynôme $Q$. Cela est clair, puisque nous avons une factorisation
$P = QR$  dans $K[x]$, et en particulier dans $E[x]$. Comme la factorisation
est unique dans $E[x]$, le polynôme $P$ n'a qu'un nombre fini de facteurs
unitaires dans $E[x]$.

\medskip\noindent
Si $F$ est un corps fini (ce qui équivaut à dire que $E$ est un corps fini), alors
$E/F$ possède un élément primitif d'après la discussion de la
Section \ref{section-finite}.
Supposons ensuite $F$ infini et qu'il n'existe qu'un nombre fini de sous-corps
propres $E/K/F$. Énumérons-les, disons $K_1, \ldots, K_N$. Alors
chaque $K_i \subset E$ est un sous-espace vectoriel propre sur $F$. Comme $F$ est infini,
nous pouvons trouver un vecteur $\alpha \in E$ tel que $\alpha \not \in K_i$ pour tout $i$
(un espace vectoriel ne peut jamais être la réunion d'un nombre fini de sous-espaces propres ;
nous omettons les détails). Alors $\alpha$ est un élément primitif de $E$ sur $F$.

\medskip\noindent
Après avoir établi l'équivalence de (1) et (2), démontrons maintenant
la dernière assertion du lemme. Choisissons une clôture algébrique
$\overline{F}$ de $F$. Énumérons les éléments
$\sigma_1, \ldots, \sigma_n \in \Mor_F(E, \overline{F})$.
Puisque $E/F$ est séparable, on a $n = [E : F]$ d'après le
Lemme \ref{lemma-separable-equality}.
Notons que, si $i \not = j$, alors
$$
V_{ij} = \Ker(\sigma_i - \sigma_j : E \longrightarrow \overline{F})
$$
n'est pas égal à $E$. En raisonnant comme au paragraphe précédent,
nous pouvons donc trouver $\alpha \in E$ tel que $\alpha \not \in V_{ij}$ pour tous
$i \not = j$. Il s'ensuit que $|\Mor_F(F(\alpha), \overline{F})| \geq n$.
D'autre part, $[F(\alpha) : F] \leq [E : F]$. Il y a donc égalité
d'après le Lemme \ref{lemma-separable-equality},
et nous concluons que $E = F(\alpha)$.
\end{proof}






\section{Trace et norme}
\label{section-trace-pairing}

\noindent
Soit $L/K$ une extension finie de corps. D'après le
Lemme \ref{lemma-vector-space-is-free},
nous pouvons choisir un isomorphisme $L \cong K^{\oplus n}$ de $K$-modules.
Bien entendu, $n = [L : K]$ est le degré de l'extension de corps.
À l'aide de cet isomorphisme, nous obtenons un morphisme de $K$-algèbres
$$
L \longrightarrow \text{Mat}(n \times n, K),\quad
\alpha \longmapsto \text{matrice de la multiplication par }\alpha
$$
Ainsi, étant donné $\alpha \in L$, nous pouvons prendre la trace et le déterminant
de la matrice correspondante. Bien entendu, ces quantités ne dépendent pas
du choix de la base effectué ci-dessus. Plus canoniquement, en considérant simplement
$L$ comme un espace vectoriel de dimension finie sur $K$, nous avons
$\text{Trace}_K(\alpha : L \to L)$ et le déterminant
$\det_K(\alpha : L \to L)$.

\begin{definition}
\label{definition-trace-norm}
Soit $L/K$ une extension finie de corps. Pour $\alpha \in L$, nous définissons
la {\it trace}
$\text{Trace}_{L/K}(\alpha) = \text{Trace}_K(\alpha : L \to L)$
et la {\it norme}
$\text{Norm}_{L/K}(\alpha) = \det_K(\alpha : L \to L)$.
\end{definition}

\noindent
Il résulte clairement de la définition que
$\text{Trace}_{L/K}$ est $K$-linéaire et vérifie
$\text{Trace}_{L/K}(\alpha) = [L : K]\alpha$ pour $\alpha \in K$.
De même, $\text{Norm}_{L/K}$ est multiplicative et
$\text{Norm}_{L/K}(\alpha) = \alpha^{[L : K]}$ pour $\alpha \in K$.
C'est un cas particulier de la construction plus générale examinée
dans les Exercices \ref{exercises-exercise-trace-det} et
\ref{exercises-exercise-trace-det-rings}.

\begin{lemma}
\label{lemma-characteristic-vs-minimal-polynomial}
Soit $L/K$ une extension finie de corps. Soit $\alpha \in L$, et soit $P$
le polynôme minimal de $\alpha$ sur $K$. Alors le polynôme caractéristique
de l'application $K$-linéaire $\alpha : L \to L$ est égal à
$P^e$, avec $e \deg(P) = [L : K]$.
\end{lemma}

\begin{proof}
Choisissons une base $\beta_1, \ldots, \beta_e$ de $L$ sur $K(\alpha)$.
Alors $e$ vérifie $e \deg(P) = [L : K]$ d'après les
Lemmes \ref{lemma-degree-minimal-polynomial} et
\ref{lemma-multiplicativity-degrees}.
Nous voyons alors que $L = \bigoplus K(\alpha) \beta_i$ est une
décomposition en somme directe de sous-espaces invariants par $\alpha$ ;
le polynôme caractéristique de $\alpha : L \to L$
est donc égal au polynôme caractéristique de
$\alpha : K(\alpha) \to K(\alpha)$ élevé à la puissance $e$.

\medskip\noindent
Pour achever la démonstration, nous pouvons supposer que $L = K(\alpha)$.
Dans ce cas, le théorème de Cayley-Hamilton montre que $\alpha$
est une racine du polynôme caractéristique. Et puisque le
polynôme caractéristique a le même degré que le polynôme
minimal, nous obtenons l'égalité.
\end{proof}

\begin{lemma}
\label{lemma-trace-and-norm-from-minimal-polynomial}
Soit $L/K$ une extension finie de corps. Soit $\alpha \in L$, et soit
$P = x^d + a_1 x^{d - 1} + \ldots + a_d$
le polynôme minimal de $\alpha$ sur $K$. Alors
$$
\text{Norm}_{L/K}(\alpha) = (-1)^{[L : K]} a_d^e
\quad\text{et}\quad
\text{Trace}_{L/K}(\alpha) = - e a_1
$$
où $e d = [L : K]$.
\end{lemma}

\begin{proof}
Cela résulte immédiatement du Lemme \ref{lemma-characteristic-vs-minimal-polynomial}
et des définitions.
\end{proof}

\begin{lemma}
\label{lemma-trace-and-norm-linear}
Soit $L/K$ une extension finie de corps. Soit $V$ un espace vectoriel de dimension
finie sur $L$. Soit $\varphi : V \to V$ une application $L$-linéaire.
Alors
$$
\text{Trace}_K(\varphi : V \to V) =
\text{Trace}_{L/K}(\text{Trace}_L(\varphi : V \to V))
$$
et
$$
\det\nolimits_K(\varphi : V \to V) =
\text{Norm}_{L/K}(\det\nolimits_L(\varphi : V \to V))
$$
\end{lemma}

\begin{proof}
Choisissons un isomorphisme $V = L^{\oplus n}$ de sorte que $\varphi$ corresponde
à une matrice $n \times n$. Dans le cas des traces, les deux membres de la formule
sont additifs en $\varphi$. Nous pouvons donc supposer que $\varphi$
corresponde à la matrice dont l'unique coefficient non nul est en position $(i, j)$.
Dans ce cas, un calcul direct montre que les deux membres sont égaux.

\medskip\noindent
Dans le cas des normes, les deux membres sont nuls si $\varphi$ possède un noyau non nul.
Nous pouvons donc supposer que $\varphi$ corresponde à un élément de
$\text{GL}_n(L)$. Les deux membres de la formule sont multiplicatifs en $\varphi$.
Puisque tout élément de $\text{GL}_n(L)$ est un produit de matrices
élémentaires, nous pouvons supposer que $\varphi$ soit de la forme
$$
E_{12}(\lambda) =
\left(
\begin{matrix}
1 & \lambda & \ldots \\
0 & 1 & \ldots \\
\ldots & \ldots & \ldots
\end{matrix}
\right)
\quad\text{ou}\quad
E_1(a) =
\left(
\begin{matrix}
a & 0 & \ldots \\
0 & 1 & \ldots \\
\ldots & \ldots & \ldots
\end{matrix}
\right)
$$
(car nous pouvons également permuter les éléments de la base à notre convenance).
Dans les deux cas, la formule se vérifie aisément par un calcul direct.
\end{proof}

\begin{lemma}
\label{lemma-trace-and-norm-tower}
Soit $M/L/K$ une tour d'extensions finies de corps. Alors
$$
\text{Trace}_{M/K} = \text{Trace}_{L/K} \circ \text{Trace}_{M/L}
\quad\text{et}\quad
\text{Norm}_{M/K} = \text{Norm}_{L/K} \circ \text{Norm}_{M/L}
$$
\end{lemma}

\begin{proof}
Considérons $M$ comme un espace vectoriel sur $L$ et appliquons le
Lemme \ref{lemma-trace-and-norm-linear}.
\end{proof}

\noindent
La forme trace se définit au moyen de la trace.

\begin{definition}
\label{definition-trace-pairing}
Soit $L/K$ une extension finie de corps. La {\it forme trace}
de $L/K$ est la forme $K$-bilinéaire symétrique
$$
Q_{L/K} : L \times L \longrightarrow K,\quad
(\alpha, \beta) \longmapsto \text{Trace}_{L/K}(\alpha\beta)
$$
\end{definition}

\noindent
Il se trouve qu'une extension finie de corps est séparable si et seulement
si la forme trace est non dégénérée.

\begin{lemma}
\label{lemma-separable-trace-pairing}
Soit $L/K$ une extension finie de corps. Les assertions suivantes sont équivalentes :
\begin{enumerate}
\item $L/K$ est séparable,
\item $\text{Trace}_{L/K}$ n'est pas identiquement nulle, et
\item la forme trace $Q_{L/K}$ est non dégénérée.
\end{enumerate}
\end{lemma}

\begin{proof}
Il est clair que (3) implique (2). Si (2) est vérifiée, choisissons $\gamma \in L$
tel que $\text{Trace}_{L/K}(\gamma) \not = 0$. Alors, si $\alpha \in L$
est non nul, nous voyons que $Q_{L/K}(\alpha, \gamma/\alpha) \not = 0$.
Ainsi, $Q_{L/K}$ est non dégénérée. Cela démontre l'équivalence de
(2) et (3).

\medskip\noindent
Supposons que $K$ soit de caractéristique $p$ et que $L = K(\alpha)$ avec
$\alpha \not \in K$ et $\alpha^p \in K$. Alors $\text{Trace}_{L/K}(1) = p = 0$.
Pour $i = 1, \ldots, p - 1$, nous voyons que $x^p - \alpha^{pi}$ est le polynôme
minimal de $\alpha^i$ sur $K$, et nous obtenons
$\text{Trace}_{L/K}(\alpha^i) = 0$ d'après le
Lemme \ref{lemma-trace-and-norm-from-minimal-polynomial}.
Ainsi, pour ce type d'extension purement inséparable de degré $p$,
nous voyons que $\text{Trace}_{L/K}$ est identiquement nulle.

\medskip\noindent
Supposons que $L/K$ ne soit pas séparable. Il existe alors un corps intermédiaire
$L/K'/K$ tel que $L/K'$ soit une extension purement inséparable de degré $p$
comme au paragraphe précédent ; voir les
Lemmes \ref{lemma-separable-first} et \ref{lemma-finite-purely-inseparable}.
Ainsi, d'après le Lemme \ref{lemma-trace-and-norm-tower},
nous voyons que $\text{Trace}_{L/K}$ est identiquement nulle.

\medskip\noindent
Supposons au contraire que $L/K$ soit séparable.
Par récurrence sur le degré, nous allons montrer que
$\text{Trace}_{L/K}$ n'est pas identiquement nulle.
Ainsi, d'après le Lemme \ref{lemma-trace-and-norm-tower}, nous pouvons supposer que
$L/K$ est engendrée par un seul élément $\alpha$ (utiliser le fait que, si
la trace est non nulle, alors elle est surjective).
Nous devons montrer que $\text{Trace}_{L/K}(\alpha^e)$
est non nulle pour un certain $e \geq 0$.
Soit $P = x^d + a_1 x^{d - 1} + \ldots + a_d$ le polynôme
minimal de $\alpha$ sur $K$.
Alors $P$ est aussi le polynôme caractéristique de l'application linéaire
$\alpha : L \to L$ ; voir le
Lemme \ref{lemma-characteristic-vs-minimal-polynomial}.
Puisque $L/K$ est séparable, le Lemme \ref{lemma-recognize-separable} montre
que $P$ a $d$ racines deux à deux distinctes $\alpha_1, \ldots, \alpha_d$
dans une clôture algébrique $\overline{K}$ de $K$. Ce sont donc les
valeurs propres de $\alpha : L \to L$.
D'après l'algèbre linéaire, la trace de $\alpha^e$ est
égale à $\alpha_1^e + \ldots + \alpha_d^e$.
Nous concluons donc par le Lemme \ref{lemma-sums-of-powers}.
\end{proof}

\noindent
Soit $K$ un corps et soit $Q : V \times V \to K$ une forme bilinéaire
sur un espace vectoriel de dimension finie sur $K$. Posons $\dim_K(V) = n$.
Alors $Q$ définit une application linéaire $Q : V \to V^*$, $v \mapsto Q(v, -)$,
où $V^* = \Hom_K(V, K)$ est l'espace vectoriel dual. D'où une application linéaire
$$
\det(Q) : \wedge^n(V) \longrightarrow \wedge^n(V)^*
$$
Si nous choisissons un vecteur de base $\omega \in \wedge^n(V)$, nous pouvons
écrire $\det(Q)(\omega) = \lambda \omega^*$, où $\omega^*$
est le vecteur de base dual dans $\wedge^n(V)^*$. Si nous remplaçons notre
choix de $\omega$ par $c \omega$ pour un certain $c \in K^*$, alors
$\omega^*$ devient $c^{-1} \omega^*$ et, par conséquent,
$\lambda$ devient $c^2 \lambda$. Ainsi, la classe de
$\lambda$ dans $K/(K^*)^2$ est bien définie et s'appelle le
{\it discriminant de $Q$}. En développant les définitions, nous voyons que
$$
\lambda = \det(Q(v_i, v_j)_{1 \leq i, j \leq n})
$$
si $\{v_1, \ldots, v_n\}$ est une base de $V$ sur $K$. Notons que
le discriminant est non nul si et seulement si $Q$ est non dégénérée.

\begin{definition}
\label{definition-discriminant}
Soit $L/K$ une extension finie de corps. Le
{\it discriminant de $L/K$} est le discriminant de
la forme trace $Q_{L/K}$.
\end{definition}

\noindent
D'après la discussion ci-dessus et le Lemme \ref{lemma-separable-trace-pairing},
nous voyons que le discriminant est non nul si et seulement
si $L/K$ est séparable. Pour $a \in K$, nous disons souvent
« le discriminant vaut $a$ » alors qu'il serait plus exact
de dire que le discriminant est la classe de $a$ dans $K/(K^*)^2$.

\begin{exercise}
\label{exercise-quadratic-discriminant}
Soit $L/K$ une extension de degré $2$. Montrer qu'il se produit exactement
l'un des cas suivants :
\begin{enumerate}
\item le discriminant vaut $0$, la caractéristique de $K$ vaut $2$,
et $L/K$ est purement inséparable et s'obtient en adjoignant une racine carrée
d'un élément de $K$,
\item le discriminant vaut $1$, la caractéristique de $K$ vaut $2$, et
$L/K$ est séparable de degré $2$,
\item le discriminant n'est pas un carré, la caractéristique de $K$
n'est pas $2$, et $L$ s'obtient à partir de $K$ en adjoignant la racine carrée
du discriminant.
\end{enumerate}
\end{exercise}







\section{Théorie de Galois}
\label{section-galois-theory}

\noindent
Voici la définition.

\begin{definition}
\label{definition-galois}
Une extension de corps $E/F$ est dite {\it galoisienne} si elle est algébrique,
séparable et normale.
\end{definition}

\noindent
On verra qu'une extension finie est galoisienne si et seulement si elle possède
le nombre « correct » d'automorphismes.

\begin{lemma}
\label{lemma-finite-Galois}
Soit $E/F$ une extension finie de corps. Alors $E$ est galoisienne sur $F$
si et seulement si $|\text{Aut}(E/F)| = [E : F]$.
\end{lemma}

\begin{proof}
Supposons que $|\text{Aut}(E/F)| = [E : F]$. D'après le
Lemme \ref{lemma-normal-and-automorphisms}, cela implique
que $E/F$ est séparable et normale, donc galoisienne.
Réciproquement, si $E/F$ est séparable, alors $[E : F] = [E : F]_s$, et
si $E/F$ est de plus normale, le
Lemme \ref{lemma-normal-and-automorphisms} implique que
$|\text{Aut}(E/F)| = [E : F]$.
\end{proof}

\noindent
Motivés par le lemme précédent, nous introduisons le groupe de Galois comme suit.

\begin{definition}
\label{definition-galois-group}
Si $E/F$ est une extension galoisienne, le groupe $\text{Aut}(E/F)$ est
appelé le {\it groupe de Galois} et est noté $\text{Gal}(E/F)$.
\end{definition}

\noindent
Si $L/K$ est une extension galoisienne infinie, il faut considérer le
groupe de Galois comme un groupe topologique. Nous y reviendrons à la
Section \ref{section-infinite-galois}.

\begin{lemma}
\label{lemma-galois-goes-up}
Soit $K/E/F$ une tour d'extensions algébriques de corps.
Si $K$ est galoisienne sur $F$, alors $K$ est galoisienne sur $E$.
\end{lemma}

\begin{proof}
Il suffit de combiner les Lemmes \ref{lemma-normal-goes-up} et \ref{lemma-separable-goes-up}.
\end{proof}

\begin{lemma}
\label{lemma-normal-closure-galois}
Soit $L/K$ une extension finie séparable de corps.
Soit $M$ la clôture normale de $L$ sur $K$
(Définition \ref{definition-normal-closure}).
Alors $M/K$ est galoisienne.
\end{lemma}

\begin{proof}
L'extension intermédiaire $M/M_{sep}/K$ du Lemme \ref{lemma-separable-first}
est normale d'après le Lemme \ref{lemma-separable-first-normal}.
Puisque $L/K$ est séparable, nous avons $L \subset M_{sep}$.
Par minimalité, $M = M_{sep}$, ce qui achève la démonstration.
\end{proof}

\noindent
Soit $G$ un groupe agissant sur un corps $K$ (par automorphismes de corps).
Nous emploierons souvent la notation
$$
K^G = \{x \in K \mid \sigma(x) = x \ \forall \sigma \in G\}
$$
et appellerons ce corps le {\it corps des invariants} de l'action de $G$ sur $K$.

\begin{lemma}
\label{lemma-galois-over-fixed-field}
Soit $K$ un corps. Soit $G$ un groupe fini agissant fidèlement sur $K$.
Alors l'extension $K/K^G$ est galoisienne, on a $[K : K^G] = |G|$,
et le groupe de Galois de l'extension est $G$.
\end{lemma}

\begin{proof}
Étant donné $\alpha \in K$, considérons l'orbite $G \cdot \alpha \subset K$
de $\alpha$ sous l'action du groupe. Considérons le polynôme
$$
P = \prod\nolimits_{\beta \in G \cdot \alpha} (x - \beta) \in K[x]
$$
Le point clé de tout le lemme est que ce polynôme est invariant
sous l'action de $G$ et possède donc ses coefficients dans $K^G$.
En effet, pour $\tau \in G$, nous avons
$$
P^\tau = \prod\nolimits_{\beta \in G \cdot \alpha} (x - \tau(\beta)) =
\prod\nolimits_{\beta \in G \cdot \alpha} (x - \beta) = P
$$
car l'application $\beta \mapsto \tau(\beta)$ est une permutation de
l'orbite $G \cdot \alpha$. Ainsi $P \in K^G[x]$. De plus,
$P(\alpha) = 0$, puisque $\alpha$ appartient à son orbite ;
nous en concluons que l'extension $K/K^G$ est algébrique. En outre,
le polynôme minimal $Q$ de $\alpha$ sur $K^G$ divise le
polynôme $P$ que nous venons de construire. Par conséquent, $Q$ est séparable
(par exemple d'après le Lemme \ref{lemma-recognize-separable}) et
nous en concluons que $K/K^G$ est séparable. Ainsi $K/K^G$ est galoisienne.
Pour achever la démonstration, il suffit de montrer que $[K : K^G] = |G|$,
car alors $G$ sera le groupe de Galois d'après le
Lemme \ref{lemma-finite-Galois}.

\medskip\noindent
Choisissons un nombre fini d'éléments $\alpha_i \in K$, $i = 1, \ldots, n$,
de sorte que les égalités $\sigma(\alpha_i) = \alpha_i$ pour $i = 1, \ldots, n$ impliquent
que $\sigma$ est l'élément neutre de $G$. Posons
$$
L = K^G(\{\sigma(\alpha_i); 1 \leq i \leq n, \sigma \in G\}) \subset K
$$
et observons que l'action de $G$ sur $K$ induit une action de $G$ sur $L$.
Nous allons montrer que $L$ est de degré $|G|$ sur $K^G$. Cela achèvera la
démonstration : en effet, si $L \subset K$ était stricte, nous pourrions ajouter un élément
$\alpha \in K$, $\alpha \not \in L$, à notre liste d'éléments
$\alpha_1, \ldots, \alpha_n$ sans agrandir $L$, ce qui est absurde.
Nous sommes ainsi ramenés au cas où $K/K^G$ est finie, lequel est
traité dans le paragraphe suivant.

\medskip\noindent
Supposons $K/K^G$ finie. D'après le Lemme \ref{lemma-primitive-element},
nous pouvons trouver $\alpha \in K$ tel que $K = K^G(\alpha)$.
La construction du premier paragraphe de cette démonstration montre
que $\alpha$ est de degré au plus $|G|$ sur $K^G$. Toutefois, ce
degré ne peut être inférieur à $|G|$, puisque $G$ agit fidèlement sur
$K^G(\alpha) = K$ par construction et en vertu de l'inégalité du
Lemme \ref{lemma-normal-and-automorphisms}.
\end{proof}

\begin{theorem}[Théorème fondamental de la théorie de Galois]
\label{theorem-galois-theory}
Soit $L/K$ une extension galoisienne finie de groupe de Galois $G$.
Alors $K = L^G$ et l'application
$$
\{\text{sous-groupes de }G\}
\longrightarrow
\{\text{extensions intermédiaires }L/M/K\},\quad
H \longmapsto L^H
$$
est une bijection dont l'inverse envoie $M$ sur $\text{Gal}(L/M)$.
Les sous-groupes normaux $H$ de $G$ correspondent exactement aux
extensions intermédiaires $M$ telles que $M/K$ soit galoisienne.
\end{theorem}

\begin{proof}
D'après le Lemme \ref{lemma-galois-goes-up}, pour toute extension intermédiaire $L/M/K$,
l'extension $L/M$ est galoisienne. Bien entendu, $L/M$ est aussi finie
(Lemme \ref{lemma-finite-goes-up}). Ainsi $|\text{Gal}(L/M)| = [L : M]$
d'après le Lemme \ref{lemma-finite-Galois}.
Réciproquement, si $H \subset G$ est un sous-groupe fini, alors
$[L : L^H] = |H|$ d'après le Lemme \ref{lemma-galois-over-fixed-field}.
Il résulte formellement de ces deux observations que nous obtenons
la correspondance bijective annoncée dans le théorème.

\medskip\noindent
Si $H \subset G$ est normal, alors $L^H$ est stable sous l'action de
$G$ et nous obtenons une application canonique $G/H \to \text{Aut}(L^H/K)$.
Cette application est nécessairement injective puisque $\text{Gal}(L/L^H) = H$. Ainsi
$|G/H| = [L^H : K]$ et $L^H$ est galoisienne d'après le
Lemme \ref{lemma-finite-Galois}.

\medskip\noindent
Réciproquement, supposons que $K \subset M \subset L$ et que $M/K$ soit galoisienne.
D'après le Lemme \ref{lemma-lift-maps}, tout
élément $\tau \in \text{Gal}(L/K)$ induit un élément
$\tau|_M \in \text{Gal}(M/K)$. On obtient ainsi un homomorphisme
de groupes de Galois $\text{Gal}(L/K) \to \text{Gal}(M/K)$ dont le
noyau est $H$. Ainsi $H$ est un sous-groupe normal.
\end{proof}

\begin{lemma}
\label{lemma-ses-galois}
Soit $L/M/K$ une tour de corps. Supposons $L/K$ et $M/K$ finies galoisiennes.
Alors nous obtenons une suite exacte courte
$$
1 \to \text{Gal}(L/M) \to \text{Gal}(L/K) \to \text{Gal}(M/K) \to 1
$$
de groupes finis.
\end{lemma}

\begin{proof}
En effet, d'après le Lemme \ref{lemma-lift-maps},
nous voyons que tout élément $\tau \in \text{Gal}(L/K)$ induit un élément
$\tau|_M \in \text{Gal}(M/K)$, ce qui fournit l'homomorphisme de
droite. L'application de gauche identifie le groupe de gauche au noyau
de la flèche de droite. La suite est exacte parce que les ordres
des groupes coïncident comme il convient, par multiplicativité des degrés
dans les tours d'extensions finies
(Lemme \ref{lemma-multiplicativity-degrees}).
On peut aussi appliquer directement le Lemme \ref{lemma-lift-maps}
pour voir que l'application de droite est surjective.
\end{proof}






\section{Théorie de Galois infinie}
\label{section-infinite-galois}

\noindent
Le groupe de Galois est muni d'une topologie canonique.

\begin{lemma}
\label{lemma-galois-profinite}
Soit $E/F$ une extension galoisienne. Munissons $\text{Gal}(E/F)$ de la topologie
la moins fine telle que
$$
\text{Gal}(E/F) \times E \longrightarrow E
$$
soit continue lorsque $E$ est muni de la topologie discrète. Alors
\begin{enumerate}
\item pour tout espace topologique $X$ et toute application $X \to \text{Gal}(E/F)$
tels que l'action $X \times E \to E$ soit continue, l'application induite
$X \to \text{Gal}(E/F)$ est continue,
\item cette topologie fait de $\text{Gal}(E/F)$
un groupe topologique profini.
\end{enumerate}
\end{lemma}

\begin{proof}
Dans toute cette démonstration, nous considérons $E$ comme un espace topologique discret.
Rappelons que la topologie compacte-ouverte sur l'ensemble des applications d'un ensemble dans lui-même
$\text{Map}(E, E)$ est la topologie universelle telle que l'action
$\text{Map}(E, E) \times E \to E$ soit continue. Voir
Topologie, Exemple \ref{topology-example-automorphisms-of-a-set}
pour un énoncé précis. La topologie du lemme sur
$\text{Gal}(E/F)$ est la topologie induite par l'application injective
$\text{Gal}(E/F) \to \text{Map}(E, E)$.
La propriété universelle (1) découle donc de la propriété universelle
correspondante de la topologie compacte-ouverte.
Comme l'ensemble des applications bijectives de l'espace discret considéré, $\text{Aut}(E)$,
muni de la topologie compacte-ouverte, forme un groupe topologique, voir
Topologie, Exemple \ref{topology-example-automorphisms-of-a-set},
et comme $\text{Gal}(E/F) = \text{Aut}(E/F) \to \text{Map}(E, E)$
se factorise par $\text{Aut}(E)$, nous obtenons un groupe topologique.
Autrement dit, nous utilisons l'injection
$$
\text{Gal}(E/F) \subset \text{Aut}(E)
$$
pour munir $\text{Gal}(E/F)$ de la structure induite de groupe topologique
(voir Topologie, Section \ref{topology-section-topological-groups})
et, par construction, il s'agit de la structure de groupe topologique
la moins fine pour laquelle l'action $\text{Gal}(E/F) \times E \to E$ est continue.

\medskip\noindent
Pour montrer que $\text{Gal}(E/F)$ est profini, nous raisonnons comme suit
(notre argument est nécessairement non standard, car nous avons défini
la topologie avant de montrer que le groupe de Galois est une limite
projective de groupes finis).
D'après Topologie, Lemme \ref{topology-lemma-profinite-group},
il suffit de montrer que l'espace topologique
sous-jacent à $\text{Gal}(E/F)$ est profini.
Pour toute partie $S \subset E$, considérons l'ensemble
$$
G(S) = \{ f : S \to E \mid
\begin{matrix}
f(\alpha)\text{ est une racine du polynôme minimal}\\
\text{de }\alpha\text{ sur }F\text{ pour tout }\alpha \in S
\end{matrix}
\}
$$
Comme un polynôme n'a qu'un nombre fini de racines, nous voyons que
$G(S)$ est fini pour toute partie finie $S \subset E$. Si $S \subset S'$,
la restriction fournit une application $G(S') \to G(S)$. Observons également
que si $\alpha \in S \cap F$ et $f \in G(S)$, alors $f(\alpha) = \alpha$
car le polynôme minimal est linéaire dans ce cas.
Considérons l'espace topologique profini
$$
G = \lim_{S \subset E\text{ fini}} G(S)
$$
Considérons l'application canonique
$$
c : \text{Gal}(E/F) \longrightarrow G,\quad
\sigma \longmapsto (\sigma|_S : S \to E)_S
$$
Elle est injective et, en explicitant les définitions, le lecteur constate
que la topologie sur $\text{Gal}(E/F)$ définie ci-dessus
est la topologie induite par $G$. Un élément $(f_S) \in G$ appartient à
l'image de $c$ exactement lorsque
(A) $f_S(\alpha) + f_S(\beta) = f_S(\alpha + \beta)$ et
(M) $f_S(\alpha)f_S(\beta) = f_S(\alpha\beta)$ chaque fois que
ces expressions ont un sens (c'est-à-dire lorsque
$\alpha, \beta, \alpha + \beta, \alpha\beta \in S$).
En effet, cela signifie que
$\lim f_S : E \to E$ sera un morphisme de $F$-algèbres,
et donc un automorphisme d'après le
Lemme \ref{lemma-algebraic-extension-self-map}.
Les conditions (A) et (M) pour un triplet donné
$(S, \alpha, \beta)$ définissent un fermé de
$G$ ; ainsi $\text{Gal}(E/F)$ est homéomorphe
à un fermé d'un espace profini, et est donc lui-même profini.
\end{proof}

\begin{lemma}
\label{lemma-galois-infinite}
Soit $L/M/K$ une tour de corps. Supposons que $L/K$ et
$M/K$ soient toutes deux galoisiennes. Alors il existe un homomorphisme canonique
continu surjectif $c : \text{Gal}(L/K) \to \text{Gal}(M/K)$.
\end{lemma}

\begin{proof}
D'après le Lemme \ref{lemma-lift-maps}, pour $\tau : L \to L$ appartenant à
$\text{Gal}(L/K)$, la restriction $\tau|_M : M \to M$
est un élément de $\text{Gal}(M/K)$. Cela définit l'homomorphisme $c$.
La continuité découle de la propriété universelle de la topologie :
l'action
$$
\text{Gal}(L/K) \times M \longrightarrow M,\quad
(\tau, x) \longmapsto \tau(x) = c(\tau)(x)
$$
est continue puisque $M \subset L$ et que l'action
$\text{Gal}(L/K) \times L \to L$ est continue.
La continuité de $c$ résulte donc de la partie (1) du
Lemme \ref{lemma-galois-profinite}.
Le Lemme \ref{lemma-lift-maps} montre également
que l'application est surjective.
\end{proof}

\noindent
Voici une manière plus classique de considérer
le groupe de Galois d'une extension galoisienne infinie.

\begin{lemma}
\label{lemma-infinite-galois-limit}
Soit $L/K$ une extension galoisienne de groupe de Galois $G$.
Soit $\Lambda$ l'ensemble des extensions intermédiaires finies galoisiennes,
c'est-à-dire que $\lambda \in \Lambda$ correspond à $L/L_\lambda/K$,
où $L_\lambda/K$ est finie galoisienne de groupe de Galois $G_\lambda$.
Définissons un ordre partiel sur $\Lambda$ par la règle
$\lambda \geq \lambda'$ si et seulement si
$L_\lambda \supset L_{\lambda'}$. Alors
\begin{enumerate}
\item $\Lambda$ est un ensemble ordonné filtrant,
\item les $L_\lambda$ forment un système d'extensions de $K$ indexé par $\Lambda$
et $L = \colim L_\lambda$,
\item les $G_\lambda$ forment un système projectif de groupes finis indexé par $\Lambda$,
les morphismes de transition sont surjectifs, et
$$
G = \lim_{\lambda \in \Lambda} G_\lambda
$$
comme groupe profini, et
\item chacune des projections $G \to G_\lambda$ est continue et surjective.
\end{enumerate}
\end{lemma}

\begin{proof}
Tout sous-corps de $L$ contenant $K$ est séparable sur $K$
(cela résulte immédiatement de la définition). Soit $S \subset L$
une partie finie. Alors $K(S)/K$ est finie et il existe
une tour $L/E/K(S)/K$ telle que $E/K$ soit finie galoisienne, voir le
Lemme \ref{lemma-normal-closure-inside-normal}.
Ainsi $E = L_\lambda$ pour un certain $\lambda \in \Lambda$.
Cela implique en particulier que l'ensemble $\Lambda$ n'est pas vide.
De plus, étant donnés $\lambda_1, \lambda_2 \in \Lambda$, nous pouvons
écrire $L_{\lambda_i} = K(S_i)$ pour des ensembles finis
$S_1, S_2 \subset L$ (Lemme \ref{lemma-finite-finitely-generated}).
Il existe alors un $\lambda \in \Lambda$ tel que
$K(S_1 \cup S_2) \subset L_\lambda$. Par conséquent,
$\lambda \geq \lambda_1, \lambda_2$ et
$\Lambda$ est filtrant (Catégories, Définition
\ref{categories-definition-directed-system}).
Enfin, puisque tout élément de $L$ appartient à
$L_\lambda$ pour un certain $\lambda \in \Lambda$, il résulte
de la description des limites inductives filtrantes dans
Catégories, Section \ref{categories-section-directed-colimits},
que $\colim L_\lambda = L$.

\medskip\noindent
Si $\lambda \geq \lambda'$ dans $\Lambda$, nous obtenons une
application canonique surjective $G_\lambda \to G_{\lambda'}$,
$\sigma \mapsto \sigma|_{L_{\lambda'}}$,
d'après le Lemme \ref{lemma-ses-galois}. Nous obtenons ainsi un système
projectif de groupes finis dont les morphismes de transition sont surjectifs.

\medskip\noindent
Rappelons que $G = \text{Aut}(L/K)$. D'après le Lemme \ref{lemma-galois-infinite},
la restriction $\sigma|_{L_\lambda}$ d'un $\sigma \in G$ à $L_\lambda$
est un élément de $G_\lambda$. De plus, ce procédé fournit une surjection
continue $G \to G_\lambda$. Puisque les morphismes de transition
du système projectif des $G_\lambda$ sont eux aussi donnés par restriction,
il est clair que nous obtenons une application continue canonique
$$
G \longrightarrow \lim_{\lambda \in \Lambda} G_\lambda
$$
La continuité résulte de la définition des limites dans la catégorie des groupes
topologiques ; rappelons que ces limites commutent avec le foncteur d'oubli
vers les catégories des ensembles et des espaces topologiques d'après
Topologie, Lemme \ref{topology-lemma-topological-group-limits}.
D'autre part, puisque $L = \colim L_\lambda$, il est clair
que tout élément de la limite projective (considérée comme un ensemble) définit un
automorphisme de $L$. L'application est donc bijective. Comme la topologie
des deux membres est profinie et qu'une application continue bijective
entre espaces profinis est un homéomorphisme
(Topologie, Lemme \ref{topology-lemma-bijective-map}), la démonstration est achevée.
\end{proof}

\begin{theorem}[Théorème fondamental de la théorie de Galois infinie]
\label{theorem-inifinite-galois-theory}
Soit $L/K$ une extension galoisienne. Soit $G = \text{Gal}(L/K)$
le groupe de Galois considéré comme un groupe topologique profini
(Lemme \ref{lemma-galois-profinite}). Alors $K = L^G$ et l'application
$$
\{\text{sous-groupes fermés de }G\}
\longrightarrow
\{\text{extensions intermédiaires }L/M/K\},\quad
H \longmapsto L^H
$$
est une bijection dont l'inverse envoie $M$ sur $\text{Gal}(L/M)$.
Les extensions intermédiaires finies $M$ correspondent exactement aux sous-groupes
ouverts $H \subset G$. Les sous-groupes fermés normaux $H$ de $G$
correspondent exactement aux extensions intermédiaires $M$ galoisiennes sur $K$.
\end{theorem}

\begin{proof}
Nous utiliserons le résultat de la théorie de Galois finie
(Théorème \ref{theorem-galois-theory})
sans autre mention.
Soit $S \subset L$ une partie finie. Il existe une tour
$L/E/K$ telle que $K(S) \subset E$ et que
$E/K$ soit finie galoisienne, voir le Lemme \ref{lemma-normal-closure-inside-normal}.
Autrement dit, $L/K$ est la réunion de ses extensions intermédiaires
finies galoisiennes.
Pour un tel $E$, d'après le Lemme \ref{lemma-galois-infinite},
l'application $\text{Gal}(L/K) \to \text{Gal}(E/K)$ est surjective
et continue, c'est-à-dire que son noyau est ouvert puisque la topologie
sur $\text{Gal}(E/K)$ est discrète.
En particulier, aucun élément de $L \setminus K$ n'est fixé par
$\text{Gal}(L/K)$, car $E^{\text{Gal}(E/K)} = K$.
Cela prouve que $L^G = K$.

\medskip\noindent
D'après le Lemme \ref{lemma-galois-goes-up}, pour toute extension intermédiaire $L/M/K$,
l'extension $L/M$ est galoisienne. Il résulte immédiatement de la définition
de la topologie sur $G$ que le sous-groupe $\text{Gal}(L/M)$ est fermé.
En appliquant ce qui précède à $L/M$, nous voyons que $L^{\text{Gal}(L/M)} = M$.

\medskip\noindent
Réciproquement, soit $H \subset G$ un sous-groupe fermé. Nous affirmons que
$H = \text{Gal}(L/L^H)$. L'inclusion $H \subset \text{Gal}(L/L^H)$
est claire. Supposons que $g \in \text{Gal}(L/L^H)$. Soit $S \subset L$
une partie finie. Nous allons montrer que le voisinage ouvert
$U_S(g) = \{g' \in G \mid g'(s) = g(s)\text{ pour tout }s \in S\}$ de $g$ rencontre $H$.
Cela implique que $g \in H$, puisque $H$ est fermé.
Soit $L/E/K$ une extension intermédiaire finie galoisienne contenant $K(S)$,
comme dans le premier paragraphe de la démonstration, et considérons l'homomorphisme
$c : \text{Gal}(L/K) \to \text{Gal}(E/K)$.
Alors $L^H \cap E = E^{c(H)}$. Puisque $g$ fixe $L^H$, il fixe
$E^{c(H)}$ et donc $c(g) \in c(H)$ par la théorie de Galois finie.
Choisissons $h \in H$ tel que $c(h) = c(g)$. Alors $h \in U_S(g)$, comme voulu.

\medskip\noindent
À ce stade, nous avons établi la correspondance entre sous-groupes fermés
et extensions intermédiaires.

\medskip\noindent
Supposons $H \subset G$ ouvert. En raisonnant comme ci-dessus, nous voyons que
$H$ contient $\text{Gal}(L/E)$ pour une extension intermédiaire finie
galoisienne $E$ assez grande, et que $L^H$ est contenu
dans $E$, donc fini sur $K$. Réciproquement, si $M$ est une extension
intermédiaire finie, alors $M$ est engendré par une partie finie $S$,
et le sous-groupe correspondant est l'ouvert $U_S(e)$,
où $e \in G$ est l'élément neutre.

\medskip\noindent
Supposons que $K \subset M \subset L$ et que $M/K$ soit galoisienne.
D'après le Lemme \ref{lemma-galois-infinite}, il existe un homomorphisme surjectif
continu de groupes de Galois
$\text{Gal}(L/K) \to \text{Gal}(M/K)$ dont le
noyau est $\text{Gal}(L/M)$. Ainsi $\text{Gal}(L/M)$ est un sous-groupe
fermé normal.

\medskip\noindent
Enfin, supposons $N \subset G$ normal et fermé. Pour toute
extension intermédiaire $L/E/K$ comme dans le premier paragraphe de la démonstration, l'image
$c(N) \subset \text{Gal}(E/K)$ est un sous-groupe normal.
Par conséquent, $L^N = \bigcup E^{c(N)}$ est une réunion d'extensions galoisiennes
de $K$ (par la théorie de Galois finie), et est donc galoisienne sur $K$.
\end{proof}

\begin{lemma}
\label{lemma-ses-infinite-galois}
Soit $L/M/K$ une tour de corps. Supposons $L/K$ et $M/K$ galoisiennes.
Alors nous obtenons une suite exacte courte
$$
1 \to \text{Gal}(L/M) \to \text{Gal}(L/K) \to \text{Gal}(M/K) \to 1
$$
de groupes topologiques profinis.
\end{lemma}

\begin{proof}
C'est une reformulation du Lemme \ref{lemma-galois-infinite}.
\end{proof}





\section{Les nombres complexes}
\label{section-complex-numbers}

\noindent
Le théorème fondamental de l'algèbre affirme que le corps des nombres complexes
est algébriquement clos. Nous en discutons brièvement dans cette section.

\medskip\noindent
La première remarque que nous souhaitons faire est qu'il faut utiliser quelques
résultats d'analyse pour le démontrer. Nous utiliserons le fait intuitivement
clair que tout polynôme de degré impair à coefficients réels possède une racine réelle.
En effet, soit
$P(x) = a_{2k + 1} x^{2k + 1} + \ldots + a_0 \in \mathbf{R}[x]$
pour un certain $k \geq 0$ et avec $a_{2k + 1} \not = 0$.
Nous pouvons supposer, et supposons, que $a_{2k + 1} > 0$. Alors, pour
$x \in \mathbf{R}$ positif et très grand, nous avons $P(x) > 0$, car le terme
$a_{2k + 1} x^{2k + 1}$ domine tous les autres termes. De même,
si $x \ll 0$, alors $P(x) < 0$ pour la même raison (et c'est ici
que nous utilisons que le degré est impair). Le théorème des valeurs
intermédiaires donne donc un $x \in \mathbf{R}$ tel que $P(x) = 0$.

\medskip\noindent
Une conséquence de ce qui précède est que $\mathbf{R}$ ne possède
aucune extension de corps non triviale de degré impair, car les éléments de telles extensions
auraient des polynômes minimaux de degré impair.

\medskip\noindent
Soit ensuite $K/\mathbf{R}$ une extension finie galoisienne de groupe de Galois $G$.
Soit $P \subset G$ un $2$-sous-groupe de Sylow. Alors $K^P/\mathbf{R}$ est une
extension de degré impair, donc $K^P = \mathbf{R}$ d'après ce qui précède,
ce qui implique à son tour
$G = P$. (Tous ces arguments reposent bien entendu sur la théorie de Galois.)
Ainsi $G$ est un $2$-groupe. Si $G$ est non trivial, nous voyons alors que
$\mathbf{C} \subset K$, puisque $\mathbf{C}$ est (à isomorphisme près) l'unique
extension de degré $2$ de $\mathbf{R}$. Si $G$ avait plus de $2$ éléments,
nous obtiendrions une extension quadratique de $\mathbf{C}$.
C'est absurde, car tout nombre complexe possède une racine carrée.

\medskip\noindent
Conclusion : $\mathbf{C}$ est algébriquement clos. En effet, sinon
nous obtiendrions une extension finie non triviale $K/\mathbf{C}$
que nous pourrions supposer normale (donc galoisienne) sur $\mathbf{R}$ d'après le
Lemme \ref{lemma-normal-closure}. Mais nous avons vu ci-dessus qu'alors
$K = \mathbf{C}$.

\begin{lemma}[Théorème fondamental de l'algèbre]
\label{lemma-C-algebraically-closed}
Le corps $\mathbf{C}$ est algébriquement clos.
\end{lemma}

\begin{proof}
Voir la discussion ci-dessus.
\end{proof}





\section{Extensions de Kummer}
\label{section-Kummer}

\noindent
Soit $K$ un corps. Soit $n \geq 2$ un entier tel que $K$ contienne
une racine primitive $n$-ième de $1$. Soit $a \in K^*$. Soit $L$ une extension
de $K$ obtenue en adjoignant une racine $b$ de l'équation $x^n = a$.
Alors $L/K$ est galoisienne. Si $G = \text{Gal}(L/K)$ est le groupe de Galois,
l'application
$$
G \longrightarrow \mu_n(K),\quad \sigma \longmapsto \sigma(b)/b
$$
est un homomorphisme injectif de groupes. En particulier, $G$ est cyclique
et son ordre divise $n$, puisqu'il s'identifie à un sous-groupe du groupe cyclique $\mu_n(K)$.
La théorie de Kummer fournit une réciproque.

\begin{lemma}[Extensions de Kummer]
\label{lemma-Kummer}
Soit $L/K$ une extension galoisienne de corps dont le groupe de Galois est
$\mathbf{Z}/n\mathbf{Z}$. Supposons de plus que la caractéristique de $K$
soit première à $n$ et que $K$ contienne une racine primitive $n$-ième de $1$.
Alors $L = K[z]$ avec $z^n \in K$.
\end{lemma}

\begin{proof}
Soit $\zeta \in K$ une racine primitive $n$-ième de $1$.
Soit $\sigma$ un générateur de $\text{Gal}(L/K)$.
Considérons $\sigma : L \to L$ comme un opérateur $K$-linéaire.
Notons que $\sigma^n - 1 = 0$ comme opérateur linéaire.
En appliquant l'indépendance linéaire des caractères
(Lemme \ref{lemma-independence-characters}), nous voyons
qu'aucun polynôme sur $K$ de degré $< n$ ne peut
annuler $\sigma$. Le polynôme minimal de $\sigma$
comme opérateur linéaire est donc $x^n - 1$. 
Puisque $\zeta$ est une racine de $x^n - 1$, l'algèbre linéaire
donne un $0 \neq z \in L$ tel que $\sigma(z) = \zeta z$.
Cet élément $z$ vérifie $z^n \in K$, car
$\sigma(z^n) = (\zeta z)^n = z^n$. De plus, les éléments
$z, \sigma(z), \ldots, \sigma^{n - 1}(z) =
z, \zeta z, \ldots \zeta^{n - 1} z$ sont deux à deux distincts,
ce qui garantit que $z$ engendre $L$ sur $K$.
Ainsi $L = K[z]$, ce qu'il fallait démontrer.
\end{proof}

\begin{lemma}
\label{lemma-adjoint-pth-root-unity}
Soit $K$ un corps, muni d'une clôture algébrique $\overline{K}$.
Soit $p$ un nombre premier distinct de la caractéristique de $K$.
Soit $\zeta \in \overline{K}$ une racine primitive $p$-ième
de $1$. Alors $K(\zeta)/K$ est une extension galoisienne dont le degré divise $p - 1$.
\end{lemma}

\begin{proof}
Le polynôme $x^p - 1$ se décompose entièrement sur
$K(\zeta)$, ses racines étant $1, \zeta, \zeta^2, \ldots, \zeta^{p - 1}$.
Ainsi, dans l'extension $K(\zeta)/K$, le corps supérieur est un corps de
décomposition ; cette extension est donc normale.
L'extension est séparable, car $x^p - 1$ est un polynôme séparable.
Elle est donc galoisienne. Tout automorphisme de $K(\zeta)$ sur $K$
envoie $\zeta$ sur $\zeta^i$ pour un certain $1 \leq i \leq p - 1$.
Le groupe de Galois est donc un sous-groupe de $(\mathbf{Z}/p\mathbf{Z})^*$.
\end{proof}

\begin{lemma}
\label{lemma-subfields-kummer}
\begin{reference}
\cite[Théorème 5.2]{Radical}
\end{reference}
Soit $K$ un corps. Soit $L/K$ une extension finie de degré $e$
engendrée par un élément $\alpha$ tel que $a = \alpha^e \in K$.
Si toute racine $e$-ième de l'unité appartenant à $L$ est contenue dans $K$, alors 
toute extension intermédiaire $L/L'/K$ est engendrée par $\alpha^d$ pour un certain $d | e$.
\end{lemma}

\begin{proof}
Observons que, pour $d | e$, le sous-corps $K(\alpha^d)$ vérifie
$[K(\alpha^d) : K] = e/d$ et $[L : K(\alpha^d)] = d$.
Soit $L/L'/K$ une extension intermédiaire. Posons $d = [L : L']$. Si $\alpha^d \in L'$,
alors $L' = K(\alpha^d)$ pour des raisons de degré.
Soit $P \in L'[x]$ le polynôme minimal de $\alpha$ sur $L'$.
Alors $P$ divise $x^e - a$ et $P$ est de degré $d$.
Écrivons
$$
x^e - a = \prod\nolimits_{i = 1, \ldots, e} (x - \zeta_i \alpha)
$$
dans un corps de décomposition de $x^e - a$ sur $L$. Les $\zeta_i$ sont des racines
$e$-ièmes de l'unité et, après renumérotation, nous avons
$$
P = \prod\nolimits_{i = 1, \ldots, d} (x - \zeta_i \alpha)
$$
Le terme constant de $P$ est égal à
$$
c = (-1)^d(\prod\nolimits_{i = 1, \ldots, d} \zeta_i) \alpha^d
$$
et appartient à $L' \subset L$. Comme $\alpha \in L$, cela implique que
$\zeta = \prod_{i = 1, \ldots, d} \zeta_i$ appartient à $L$, et donc à $K$ par
notre hypothèse. Ainsi $\alpha^d = (-1)^d\zeta^{-1}c \in L'$, ce qui conclut.
\end{proof}




\section{Extensions d’Artin–Schreier}
\label{section-Artin-Schreier}

\noindent
Soit $K$ un corps de caractéristique $p > 0$. Soit $a \in K$. Soit $L$ une
extension de $K$ obtenue en adjoignant une racine $b$ de l'équation
$x^p - x = a$. Alors $L/K$ est galoisienne. Si $G = \text{Gal}(L/K)$ est le groupe
de Galois, alors l'application
$$
G \longrightarrow \mathbf{Z}/p\mathbf{Z},\quad
\sigma \longmapsto \sigma(b) - b
$$
est un homomorphisme injectif de groupes. En particulier, $G$ est cyclique
d'ordre divisant $p$ en tant que sous-groupe de $\mathbf{Z}/p\mathbf{Z}$.
La théorie des extensions d'Artin–Schreier fournit une réciproque.

\begin{lemma}[Extensions d'Artin–Schreier]
\label{lemma-Artin-Schreier}
Soit $L/K$ une extension galoisienne de corps de caractéristique $p > 0$
de groupe de Galois $\mathbf{Z}/p\mathbf{Z}$. Alors $L = K[z]$ avec
$z^p - z \in K$.
\end{lemma}

\begin{proof}
Soit $\sigma$ un générateur de $\text{Gal}(L/K)$.
Considérons $\sigma : L \to L$ comme un opérateur $K$-linéaire.
Observons que $\sigma^p - 1 = 0$ comme opérateur linéaire.
En appliquant l'indépendance linéaire des caractères
(Lemme \ref{lemma-independence-characters}),
aucun polynôme de degré $< p$ ne peut
annuler $\sigma$. Nous en concluons que le polynôme minimal
de $\sigma$ est $x^p - 1 = (x - 1)^p$.
Cela implique qu'il existe $w \in L$ tel que
$(\sigma - 1)^{p - 1}(w) = y$ soit non nul. Alors
$\sigma(y) = y$, c'est-à-dire $y \in K$. Ainsi
$z = y^{-1}(\sigma - 1)^{p - 2}(w)$ vérifie
$\sigma(z) = z + 1$. Comme $z \not \in K$, nous avons $L = K[z]$.
De plus, puisque $\sigma(z^p - z) = (z + 1)^p - (z + 1) = z^p - z$,
nous voyons que $z^p - z \in K$, ce qui achève la démonstration.
\end{proof}






\section{Transcendance}
\label{section-transcendence}

\noindent
Rappelons les définitions usuelles.

\begin{definition}
\label{definition-transcendence}
Soit $K/k$ une extension de corps.
\begin{enumerate}
\item Une famille d'éléments $\{x_i\}_{i \in I}$ de $K$ est dite
{\it algébriquement indépendante} sur $k$ si l'application
$$
k[X_i; i\in I] \longrightarrow K
$$
qui envoie $X_i$ sur $x_i$ est injective.
\item Le corps des fractions d'un anneau de polynômes
$k[x_i; i \in I]$ est noté $k(x_i; i\in I)$.
\item Une {\it extension transcendante pure} de $k$ est toute
extension de corps $K/k$ isomorphe au corps des
fractions d'un anneau de polynômes sur $k$.
\item Une {\it base de transcendance} de $K/k$ est une
famille d'éléments $\{x_i\}_{i \in I}$ qui sont
algébriquement indépendants sur $k$ et tels que
l'extension $K/k(x_i; i\in I)$ soit algébrique.
\end{enumerate}
\end{definition}

\begin{example}
\label{example-pi-transcendental}
Le corps $\mathbf{Q}(\pi)$ est une extension transcendante pure, car
$\pi$ n'est racine d'aucun polynôme non nul à coefficients rationnels.
En particulier, $\mathbf{Q}(\pi) \cong \mathbf{Q}(x)$.
\end{example}

\begin{lemma}
\label{lemma-transcendence-degree}
Soit $E/F$ une extension de corps. Il existe une base de transcendance de $E$ sur $F$.
Deux bases de transcendance quelconques ont même cardinal.
\end{lemma}

\begin{proof}
Soit $A$ une partie algébriquement indépendante de $E$. Soit $G$ une partie
de $E$ contenant $A$ qui engendre $E/F$. Nous affirmons qu'il existe une
base de transcendance $B$ telle que $A \subset B \subset G$.
Pour le démontrer, considérons l'ensemble $\mathcal{B}$ des parties
algébriquement indépendantes qui sont des parties de $G$ contenant $A$.
Ordonnons partiellement $\mathcal{B}$ par inclusion.
Alors $\mathcal{B}$ contient au moins l'élément $A$.
La réunion des éléments d'une partie totalement ordonnée $T$ de $\mathcal{B}$
est une partie de $E$ algébriquement indépendante sur $F$, car toute relation
de dépendance algébrique apparaîtrait déjà dans l'un des éléments de $T$
(puisque les polynômes ne font intervenir qu'un nombre fini de variables). Cette réunion
contient aussi $A$ et est contenue dans $G$. D'après le lemme de Zorn, il existe un
élément maximal $B \in \mathcal{B}$. Nous affirmons maintenant que $E$ est algébrique sur $F(B)$.
En effet, sinon il existerait un élément
$f \in G$ transcendant sur $F(B)$, puisque $F(G) = E$. Alors
$B \cup\{f\}$ serait algébriquement indépendante, contrairement à la
maximalité de $B$. Ainsi $B$ est la base de transcendance recherchée.

\medskip\noindent
Soient $B$ et $B'$ deux bases de transcendance. Sans perte de généralité, nous
pouvons supposer que $|B'| \leq |B|$. Divisons la démonstration en deux cas :
le premier est celui où $B$ est un ensemble infini. Alors, pour chaque $\alpha \in B'$,
il existe un ensemble fini $B_{\alpha} \subset B$ tel que $\alpha$ soit algébrique
sur
$F(B_{\alpha})$, puisque toute relation de dépendance algébrique n'utilise qu'un nombre fini
d'indéterminées. Posons alors $B^* = \bigcup_{\alpha\in B'} B_{\alpha}$.
Par construction, $B^* \subset B$, mais nous affirmons qu'en fait ces deux ensembles sont
égaux. Pour le voir, supposons qu'ils ne le soient pas et qu'il existe, par exemple, un élément
$\beta \in B \setminus B^*$. Nous savons que $\beta$ est algébrique sur $F(B')$, qui
est algébrique sur $F(B^*)$. Par conséquent, $\beta$ est algébrique sur $F(B^*)$, ce qui est
une contradiction. Ainsi $|B| \leq |\bigcup_{\alpha \in B'} B_{\alpha}|$.
Si $B'$ est fini, alors $B$ l'est également ; nous pouvons donc supposer $B'$ infini.
Il vient
$$
|B| \leq |\bigcup\nolimits_{\alpha \in B'} B_{\alpha}| = |B'|
$$
car chaque $B_\alpha$ est fini et $B'$ est infini. Par conséquent, dans le
cas infini, $|B| = |B'|$.

\medskip\noindent
Il reste à étudier le cas où $B$ est fini.
Dans ce cas, $B'$ est également fini ; écrivons
$B = \{\alpha_1, \ldots, \alpha_n\}$ et
$B' = \{\beta_1, \ldots, \beta_m\}$ avec $m \leq n$.
Raisonnons par récurrence sur $m$ : si $m = 0$, alors $E/F$ est algébrique, donc
$B = \emptyset$ et $n = 0$. Si $m > 0$, il existe un polynôme irréductible
$f \in F[x, y_1, \ldots, y_n]$ tel que
$f(\beta_1, \alpha_1, \ldots, \alpha_n) = 0$ et tel que $x$ intervienne dans $f$.
Puisque $\beta_1$ n'est pas algébrique sur $F$, $f$ doit faire intervenir un certain $y_i$ ;
sans perte de généralité, supposons que $f$ fasse intervenir $y_1$.
Posons $B^* = \{\beta_1, \alpha_2, \ldots, \alpha_n\}$.
Nous affirmons que $B^*$ est une base de transcendance de l'extension $E/F$. Pour le montrer, remarquons que
nous avons une tour d'extensions algébriques
$$
E/ F(B^*, \alpha_1) / F(B^*)
$$
puisque $\alpha_1$ est algébrique sur $F(B^*)$.
Nous affirmons maintenant que $B^*$ (en comptant les éléments avec multiplicité) est
algébriquement indépendante sur $F$ : en effet, sinon, il existerait un
polynôme irréductible $g\in F[x, y_2, \ldots, y_n]$ tel que
$g(\beta_1, \alpha_2, \ldots, \alpha_n) = 0$,
dans lequel $x$ devrait intervenir, rendant $\beta_1$
algébrique sur $F(\alpha_2, \ldots, \alpha_n)$, ce qui rendrait $\alpha_1$
algébrique sur $F(\alpha_2, \ldots, \alpha_n)$, chose impossible.
Cela signifie donc que $\{\alpha_2, \ldots, \alpha_n\}$ et
$\{\beta_2, \ldots, \beta_m\}$ sont des bases de transcendance de $E$ sur $F(\beta_1)$ ;
par récurrence, $m = n$.
\end{proof}

\begin{definition}
\label{definition-transcendence-degree}
Soit $K/k$ une extension de corps.
Le {\it degré de transcendance} de $K$ sur $k$ est
le cardinal d'une base de transcendance de $K$ sur $k$.
Il est noté $\text{trdeg}_k(K)$.
\end{definition}

\begin{lemma}
\label{lemma-transcendence-degree-tower}
Soient $L/K/k$ des extensions de corps.
Alors
$$
\text{trdeg}_k(L) =
\text{trdeg}_K(L) +
\text{trdeg}_k(K).
$$
\end{lemma}

\begin{proof}
Choisissons une base de transcendance $A \subset K$ de $K$ sur $k$.
Choisissons une base de transcendance $B \subset L$ de $L$ sur $K$.
Il est alors immédiat que $A \cup B$ est une base de transcendance
de $L$ sur $k$.
\end{proof}

\begin{example}
\label{example-pi-e-transcendental}
Considérons l'extension de corps $\mathbf{Q}(e, \pi)$ obtenue en
adjoignant les nombres $e$ et $\pi$. Cette extension de corps a un degré de transcendance
au moins égal à $1$, puisque $e$ et $\pi$ sont tous deux transcendants sur les
rationnels. Toutefois, cette extension de corps pourrait avoir un degré de transcendance
égal à $2$ si $e$ et $\pi$ sont algébriquement indépendants. On ignore si
c'est le cas ; le problème de la détermination de
$\text{trdeg}(\mathbf{Q}(e, \pi))$ est donc ouvert.
\end{example}

\begin{example}
\label{example-function-field-not-unique-transcendence-basis}
Soit $F$ un corps et soit $E = F(t)$. Alors $\{t\}$ est une
base de transcendance puisque $E = F(t)$. Toutefois, $\{t^2\}$
est aussi une base de transcendance puisque $F(t)/F(t^2)$ est algébrique.
Cela montre que, bien que nous puissions toujours décomposer une extension
$E/F$ en une extension algébrique $E/F'$ et une
extension transcendante pure $F'/F$, cette décomposition n'est pas unique et
dépend du choix d'une base de transcendance.
\end{example}

\begin{example}
\label{example-riemann-surface-transcendence}
Soit $X$ une surface de Riemann compacte. Alors le corps des fonctions
$\mathbf{C}(X)$ (voir Exemple \ref{example-field-of-meromorphic-functions})
est de degré de transcendance un sur $\mathbf{C}$. En fait, {\it toute}
extension de type fini de $\mathbf{C}$ de degré de transcendance
un provient d'une surface de Riemann. Il existe même une équivalence de
catégories entre la catégorie des surfaces de Riemann compactes et des
applications holomorphes (non constantes), et l'opposée de la catégorie des extensions
de type fini de $\mathbf{C}$ de degré de transcendance $1$
et des morphismes de $\mathbf{C}$-algèbres. Voir \cite{Forster}.

\medskip\noindent
Il existe également une version algébrique de l'énoncé précédent. Étant donnée une
courbe algébrique (irréductible) dans l'espace projectif sur un corps algébriquement
clos $k$ (par exemple le corps des nombres complexes), on peut considérer son
« corps des fonctions rationnelles » : il s'agit essentiellement des fonctions qui se présentent comme
des quotients de polynômes, dont le dénominateur ne s'annule pas identiquement
sur la courbe. Il existe une anti-équivalence analogue de catégories
(Courbes algébriques, Théorème \ref{curves-theorem-curves-rational-maps})
entre les courbes projectives lisses et
les morphismes non constants de courbes, et les extensions de type fini de $k$ de
degré de transcendance un. Voir \cite{H}.
\end{example}

\begin{definition}
\label{definition-algebraically-closed-in}
Soit $K/k$ une extension de corps.
\begin{enumerate}
\item La {\it clôture algébrique de $k$ dans $K$} est le sous-corps
$k'$ de $K$ formé des éléments de $K$ qui sont algébriques sur $k$.
\item Nous disons que $k$ est {\it algébriquement clos dans $K$} si
tout élément de $K$ qui est algébrique sur $k$
appartient à $k$.
\end{enumerate}
\end{definition}

\begin{lemma}
\label{lemma-purely-transcendental-degree}
Soit $k'/k$ une extension finie de corps. Soit
$k'(x_1, \ldots, x_r)/k(x_1, \ldots, x_r)$
l'extension induite des extensions transcendantes pures.
Alors $[k'(x_1, \ldots, x_r) : k(x_1, \ldots, x_r)] = [k' : k] < \infty$.
\end{lemma}

\begin{proof}
Par multiplicativité des degrés des extensions
(Lemme \ref{lemma-multiplicativity-degrees}),
il suffit de le démontrer lorsque $k'$ est engendré par un unique élément
$\alpha \in k'$ sur $k$. Soit $f \in k[T]$ le polynôme minimal
de $\alpha$ sur $k$. Alors $k'(x_1, \ldots, x_r)$ est engendré
par $\alpha, x_1, \ldots, x_r$ sur $k$, et donc $k'(x_1, \ldots, x_r)$
est engendré par $\alpha$ sur $k(x_1, \ldots, x_r)$.
Il suffit donc de montrer que $f$ reste irréductible comme
élément de $k(x_1, \ldots, x_r)[T]$. Nous ne faisons qu'esquisser la démonstration.
Il est clair que $f$ est irréductible comme élément de
$k[x_1, \ldots, x_r, T]$, par exemple parce que $f$ est unitaire
comme polynôme en $T$ et que toute factorisation éventuelle dans
$k[x_1, \ldots, x_r, T]$ conduirait à une factorisation dans $k[T]$
en posant les $x_i$ égaux à $0$. Le lemme de Gauss permet de conclure.
\end{proof}

\begin{lemma}
\label{lemma-algebraic-closure-in-finitely-generated}
Soit $K/k$ une extension de corps de type fini.
La clôture algébrique de $k$ dans $K$ est finie sur $k$.
\end{lemma}

\begin{proof}
Supposons que $x_1, \ldots, x_r \in K$ forment une base de transcendance de $K$
sur $k$. Alors $n = [K : k(x_1, \ldots, x_r)] < \infty$.
Supposons que $k \subset k' \subset K$ et que $k'/k$ soit finie.
Dans ce cas,
$[k'(x_1, \ldots, x_r) : k(x_1, \ldots, x_r)] = [k' : k] < \infty$, voir le
Lemme \ref{lemma-purely-transcendental-degree}.
Par conséquent,
$$
[k' : k] = [k'(x_1, \ldots, x_r) : k(x_1, \ldots, x_r)] \leq
[K : k(x_1, \ldots, x_r)] = n.
$$
Autrement dit, les degrés des extensions intermédiaires finies sont bornés,
et le lemme en résulte.
\end{proof}






\section{Extensions linéairement disjointes}
\label{section-linearly-disjoint}

\noindent
Soit $k$ un corps, et soient $K$ et $L$ des extensions de corps de $k$.
Supposons aussi que $K$ et $L$ soient plongés dans un corps plus grand $\Omega$.

\begin{definition}
\label{definition-compositum}
Considérons un diagramme
\begin{equation}
\label{equation-inside-omega}
\vcenter{
\xymatrix{
L \ar[r] & \Omega \\
k \ar[r] \ar[u] & K \ar[u]
}
}
\end{equation}
d'extensions de corps. Le {\it composé de $K$ et $L$ dans $\Omega$},
noté $KL$, est le plus petit sous-corps de $\Omega$ contenant à la fois
$L$ et $K$.
\end{definition}

\noindent
Il est clair que $KL$ est engendré par l'ensemble $K \cup L$ sur $k$,
par l'ensemble $K$ sur $L$, et par l'ensemble $L$ sur $K$.

\medskip\noindent
{\bf Avertissement :} La classe d'isomorphisme du corps composé dépend du
choix des plongements de $K$ et $L$ dans $\Omega$. Considérons par exemple
les corps de nombres $K = \mathbf{Q}(2^{1/8}) \subset \mathbf{R}$ et
$L = \mathbf{Q}(2^{1/12}) \subset \mathbf{R}$. Leur composé dans
$\mathbf{R}$ est le corps $\mathbf{Q}(2^{1/24})$, de degré $24$ sur
$\mathbf{Q}$. Toutefois, si nous plongeons $K = \mathbf{Q}[x]/(x^8 - 2)$ dans
$\mathbf{C}$ en envoyant $x$ sur $2^{1/8}e^{2\pi i/8}$, alors le composé
$\mathbf{Q}(2^{1/12}, 2^{1/8}e^{2\pi i/8})$ contient $i = e^{2\pi i/4}$ et est de
degré $48$ sur $\mathbf{Q}$ (nous omettons de montrer que le degré vaut $48$, mais
l'existence de $i$ prouve assurément que les deux composés ne sont pas isomorphes).

\begin{definition}
\label{definition-linearly-disjoint}
Considérons un diagramme de corps comme dans (\ref{equation-inside-omega}).
Nous disons que $K$ et $L$ sont {\it linéairement disjoints sur $k$ dans $\Omega$}
si l'application
$$
K \otimes_k L \longrightarrow KL,\quad
\sum x_i \otimes y_i \longmapsto \sum x_i y_i
$$
est injective.
\end{definition}

\noindent
Le lemme suivant ne semble trouver sa place nulle part ailleurs.

\begin{lemma}
\label{lemma-normal-case}
Soit $E/F$ une extension algébrique normale de corps. Il existe des extensions intermédiaires
$E / E_{sep} /F$ et $E / E_{insep} / F$ telles que
\begin{enumerate}
\item $F \subset E_{sep}$ soit galoisienne et $E_{sep} \subset E$
soit purement inséparable,
\item $F \subset E_{insep}$ soit purement inséparable et $E_{insep} \subset E$
soit galoisienne,
\item l'application de multiplication $E_{sep} \otimes_F E_{insep} \longrightarrow E$
est un isomorphisme.
\end{enumerate}
\end{lemma}

\begin{proof}
Nous avons trouvé le sous-corps $E_{sep}$ dans le Lemme \ref{lemma-separable-first}.
Nous posons $E_{insep} = E^{\text{Aut}(E/F)}$. Les détails sont omis.
\end{proof}









\section{Rappels}
\label{section-algebraic}

\noindent
Dans cette section, nous récapitulons brièvement les résultats obtenus ci-dessus.

\medskip\noindent
Soit $K/k$ une extension de corps. Soit $\alpha \in K$. Les
possibilités suivantes se présentent :
\begin{enumerate}
\item L'élément $\alpha$ est transcendant sur $k$.
\item L'élément $\alpha$ est algébrique sur $k$. Notons
$P(T) \in k[T]$ son {\it polynôme minimal}. C'est un polynôme unitaire
$P(T) = T^d + a_1 T^{d - 1} + \ldots + a_d$ à coefficients dans
$k$. Il est irréductible et $P(\alpha) = 0$. Ces propriétés
déterminent $P$ de manière unique, et l'entier $d$ est appelé le
{\it degré de $\alpha$ sur $k$}. Il existe deux sous-cas :
\begin{enumerate}
\item Le polynôme $\text{d}P/\text{d}T$ n'est pas identiquement nul.
Cela équivaut à la condition
$P(T) = \prod_{i = 1, \ldots, d} (T - \alpha_i)$ pour des
éléments deux à deux distincts $\alpha_1, \ldots, \alpha_d$
de la clôture algébrique de $k$.
Dans ce cas, nous disons que $\alpha$ est {\it séparable} sur $k$.
\item Le $\text{d}P/\text{d}T$ est identiquement nul. Dans ce cas, la
caractéristique $p$ de $k$ est $ > 0$, et $P$ est en fait un polynôme
en $T^p$. Il existe manifestement une plus grande puissance $q = p^e$ telle que $P$ soit
un polynôme en $T^q$. Alors l'élément $\alpha^q$ est séparable sur $k$.
\end{enumerate}
\end{enumerate}

\begin{definition}
\label{definition-separable-algebraic}
Extensions algébriques de corps.
\begin{enumerate}
\item Une extension de corps $K/k$ est dite {\it algébrique}
si tout élément de $K$ est algébrique sur $k$.
\item Une extension algébrique $k'/k$ est dite {\it séparable}
si tout $\alpha \in k'$ est séparable sur $k$.
\item Une extension algébrique
$k'/k$ est dite {\it purement inséparable} si
la caractéristique de $k$ est $p > 0$ et si, pour tout élément
$\alpha \in k'$, il existe une puissance $q$ de $p$ telle que
$\alpha^q \in k$.
\item Une extension algébrique $k'/k$ est dite {\it normale}
si, pour tout $\alpha \in k'$, le polynôme minimal $P(T) \in k[T]$
de $\alpha$ sur $k$ se décompose entièrement en facteurs linéaires sur $k'$.
\item Une extension algébrique $k'/k$ est dite {\it galoisienne}
si elle est séparable et normale.
\end{enumerate}
\end{definition}

\noindent
Le lemme suivant ne semble trouver sa place nulle part ailleurs.

\begin{lemma}
\label{lemma-pth-root}
Soit $K$ un corps de caractéristique $p > 0$. Soit $L/K$ une extension
algébrique séparable. Soit $\alpha \in L$.
\begin{enumerate}
\item Si les coefficients du polynôme minimal de $\alpha$
sur $K$ sont des puissances $p$-ièmes dans $K$, alors $\alpha$ est une puissance
$p$-ième dans $L$.
\item Plus généralement, si $P \in K[T]$ est un polynôme tel que (a) $\alpha$
soit une racine de $P$, (b) $P$ possède des racines deux à deux distinctes dans une clôture algébrique,
et (c) tous les coefficients de $P$ soient des puissances $p$-ièmes, alors $\alpha$ est une
puissance $p$-ième dans $L$.
\end{enumerate}
\end{lemma}

\begin{proof}
Il résulte des définitions que (2) implique (1). Supposons que $P$ satisfasse aux hypothèses de (2).
Écrivons $P(T) = \sum\nolimits_{i = 0}^d a_i T^{d - i}$ et $a_i = b_i^p$.
Le polynôme $Q(T) = \sum\nolimits_{i = 0}^d b_i T^{d - i}$ possède lui aussi des
racines distinctes dans une clôture algébrique, car les racines de $Q$
sont les racines $p$-ièmes des racines de $P$. Si $\alpha$ n'est pas une puissance $p$-ième,
alors $T^p - \alpha$ est un polynôme irréductible sur $L$
(Lemme \ref{lemma-take-pth-root}).
De plus, $Q$ et $T^p - \alpha$ ont une racine commune dans
une clôture algébrique $\overline{L}$.
Ainsi $Q$ et $T^p - \alpha$ ne sont pas premiers entre eux, ce qui
implique que $T^p - \alpha | Q$ dans $L[T]$. Cela contredit le fait que
les racines de $Q$ sont deux à deux distinctes.
\end{proof}











\begin{multicols}{2}[\section{Autres chapitres}]
\noindent
Préliminaires
\begin{enumerate}
\item \hyperref[introduction-section-phantom]{Introduction}
\item \hyperref[conventions-section-phantom]{Conventions}
\item \hyperref[sets-section-phantom]{Théorie des ensembles}
\item \hyperref[categories-section-phantom]{Catégories}
\item \hyperref[topology-section-phantom]{Topologie}
\item \hyperref[sheaves-section-phantom]{Faisceaux sur les espaces}
\item \hyperref[sites-section-phantom]{Sites et faisceaux}
\item \hyperref[stacks-section-phantom]{Champs}
\item \hyperref[fields-section-phantom]{Corps}
\item \hyperref[algebra-section-phantom]{Algèbre commutative}
\item \hyperref[brauer-section-phantom]{Groupes de Brauer}
\item \hyperref[homology-section-phantom]{Algèbre homologique}
\item \hyperref[derived-section-phantom]{Catégories dérivées}
\item \hyperref[simplicial-section-phantom]{Méthodes simpliciales}
\item \hyperref[more-algebra-section-phantom]{Compléments d'algèbre}
\item \hyperref[smoothing-section-phantom]{Lissage des morphismes d'anneaux}
\item \hyperref[modules-section-phantom]{Faisceaux de modules}
\item \hyperref[sites-modules-section-phantom]{Modules sur les sites}
\item \hyperref[injectives-section-phantom]{Objets injectifs}
\item \hyperref[cohomology-section-phantom]{Cohomologie des faisceaux}
\item \hyperref[sites-cohomology-section-phantom]{Cohomologie sur les sites}
\item \hyperref[dga-section-phantom]{Algèbre différentielle graduée}
\item \hyperref[dpa-section-phantom]{Algèbre à puissances divisées}
\item \hyperref[sdga-section-phantom]{Faisceaux différentiels gradués}
\item \hyperref[hypercovering-section-phantom]{Hyperrecouvrements}
\end{enumerate}
Schémas
\begin{enumerate}
\setcounter{enumi}{25}
\item \hyperref[schemes-section-phantom]{Schémas}
\item \hyperref[constructions-section-phantom]{Constructions de schémas}
\item \hyperref[properties-section-phantom]{Propriétés des schémas}
\item \hyperref[morphisms-section-phantom]{Morphismes de schémas}
\item \hyperref[coherent-section-phantom]{Cohomologie des schémas}
\item \hyperref[divisors-section-phantom]{Diviseurs}
\item \hyperref[limits-section-phantom]{Limites de schémas}
\item \hyperref[varieties-section-phantom]{Variétés}
\item \hyperref[topologies-section-phantom]{Topologies sur les schémas}
\item \hyperref[descent-section-phantom]{Descente}
\item \hyperref[perfect-section-phantom]{Catégories dérivées de schémas}
\item \hyperref[more-morphisms-section-phantom]{Compléments sur les morphismes}
\item \hyperref[flat-section-phantom]{Compléments sur la platitude}
\item \hyperref[groupoids-section-phantom]{Schémas en groupoïdes}
\item \hyperref[more-groupoids-section-phantom]{Compléments sur les schémas en groupoïdes}
\item \hyperref[etale-section-phantom]{Morphismes étales de schémas}
\end{enumerate}
Thèmes de théorie des schémas
\begin{enumerate}
\setcounter{enumi}{41}
\item \hyperref[chow-section-phantom]{Homologie de Chow}
\item \hyperref[intersection-section-phantom]{Théorie de l'intersection}
\item \hyperref[pic-section-phantom]{Schémas de Picard des courbes}
\item \hyperref[weil-section-phantom]{Théories cohomologiques de Weil}
\item \hyperref[adequate-section-phantom]{Modules adéquats}
\item \hyperref[dualizing-section-phantom]{Complexes dualisants}
\item \hyperref[duality-section-phantom]{Dualité pour les schémas}
\item \hyperref[discriminant-section-phantom]{Discriminants et différentes}
\item \hyperref[derham-section-phantom]{Cohomologie de de Rham}
\item \hyperref[local-cohomology-section-phantom]{Cohomologie locale}
\item \hyperref[algebraization-section-phantom]{Géométrie algébrique et formelle}
\item \hyperref[curves-section-phantom]{Courbes algébriques}
\item \hyperref[resolve-section-phantom]{Résolution des surfaces}
\item \hyperref[models-section-phantom]{Réduction semi-stable}
\item \hyperref[functors-section-phantom]{Foncteurs et morphismes}
\item \hyperref[equiv-section-phantom]{Catégories dérivées de variétés}
\item \hyperref[pione-section-phantom]{Groupes fondamentaux de schémas}
\item \hyperref[etale-cohomology-section-phantom]{Cohomologie étale}
\item \hyperref[crystalline-section-phantom]{Cohomologie cristalline}
\item \hyperref[proetale-section-phantom]{Cohomologie pro-étale}
\item \hyperref[relative-cycles-section-phantom]{Cycles relatifs}
\item \hyperref[more-etale-section-phantom]{Compléments de cohomologie étale}
\item \hyperref[trace-section-phantom]{La formule des traces}
\end{enumerate}
Espaces algébriques
\begin{enumerate}
\setcounter{enumi}{64}
\item \hyperref[spaces-section-phantom]{Espaces algébriques}
\item \hyperref[spaces-properties-section-phantom]{Propriétés des espaces algébriques}
\item \hyperref[spaces-morphisms-section-phantom]{Morphismes d'espaces algébriques}
\item \hyperref[decent-spaces-section-phantom]{Espaces algébriques décents}
\item \hyperref[spaces-cohomology-section-phantom]{Cohomologie des espaces algébriques}
\item \hyperref[spaces-limits-section-phantom]{Limites d'espaces algébriques}
\item \hyperref[spaces-divisors-section-phantom]{Diviseurs sur les espaces algébriques}
\item \hyperref[spaces-over-fields-section-phantom]{Espaces algébriques sur des corps}
\item \hyperref[spaces-topologies-section-phantom]{Topologies sur les espaces algébriques}
\item \hyperref[spaces-descent-section-phantom]{Descente et espaces algébriques}
\item \hyperref[spaces-perfect-section-phantom]{Catégories dérivées d'espaces}
\item \hyperref[spaces-more-morphisms-section-phantom]{Compléments sur les morphismes d'espaces}
\item \hyperref[spaces-flat-section-phantom]{Platitude sur les espaces algébriques}
\item \hyperref[spaces-groupoids-section-phantom]{Groupoïdes dans les espaces algébriques}
\item \hyperref[spaces-more-groupoids-section-phantom]{Compléments sur les groupoïdes dans les espaces}
\item \hyperref[bootstrap-section-phantom]{Amorçage}
\item \hyperref[spaces-pushouts-section-phantom]{Sommes amalgamées d'espaces algébriques}
\end{enumerate}
Thèmes de géométrie
\begin{enumerate}
\setcounter{enumi}{81}
\item \hyperref[spaces-chow-section-phantom]{Groupes de Chow des espaces}
\item \hyperref[groupoids-quotients-section-phantom]{Quotients de groupoïdes}
\item \hyperref[spaces-more-cohomology-section-phantom]{Compléments sur la cohomologie des espaces}
\item \hyperref[spaces-simplicial-section-phantom]{Espaces simpliciaux}
\item \hyperref[spaces-duality-section-phantom]{Dualité pour les espaces}
\item \hyperref[formal-spaces-section-phantom]{Espaces algébriques formels}
\item \hyperref[restricted-section-phantom]{Algébrisation des espaces formels}
\item \hyperref[spaces-resolve-section-phantom]{Résolution des surfaces revisitée}
\end{enumerate}
Théorie des déformations
\begin{enumerate}
\setcounter{enumi}{89}
\item \hyperref[formal-defos-section-phantom]{Théorie formelle des déformations}
\item \hyperref[defos-section-phantom]{Théorie des déformations}
\item \hyperref[cotangent-section-phantom]{Le complexe cotangent}
\item \hyperref[examples-defos-section-phantom]{Problèmes de déformation}
\end{enumerate}
Champs algébriques
\begin{enumerate}
\setcounter{enumi}{93}
\item \hyperref[algebraic-section-phantom]{Champs algébriques}
\item \hyperref[examples-stacks-section-phantom]{Exemples de champs}
\item \hyperref[stacks-sheaves-section-phantom]{Faisceaux sur les champs algébriques}
\item \hyperref[criteria-section-phantom]{Critères de représentabilité}
\item \hyperref[artin-section-phantom]{Axiomes d'Artin}
\item \hyperref[quot-section-phantom]{Espaces de Quot et de Hilbert}
\item \hyperref[stacks-properties-section-phantom]{Propriétés des champs algébriques}
\item \hyperref[stacks-morphisms-section-phantom]{Morphismes de champs algébriques}
\item \hyperref[stacks-limits-section-phantom]{Limites de champs algébriques}
\item \hyperref[stacks-cohomology-section-phantom]{Cohomologie des champs algébriques}
\item \hyperref[stacks-perfect-section-phantom]{Catégories dérivées de champs}
\item \hyperref[stacks-introduction-section-phantom]{Introduction aux champs algébriques}
\item \hyperref[stacks-more-morphisms-section-phantom]{Compléments sur les morphismes de champs}
\item \hyperref[stacks-geometry-section-phantom]{La géométrie des champs}
\end{enumerate}
Thèmes de théorie des espaces de modules
\begin{enumerate}
\setcounter{enumi}{107}
\item \hyperref[moduli-section-phantom]{Champs de modules}
\item \hyperref[moduli-curves-section-phantom]{Espaces de modules de courbes}
\end{enumerate}
Divers
\begin{enumerate}
\setcounter{enumi}{109}
\item \hyperref[examples-section-phantom]{Exemples}
\item \hyperref[exercises-section-phantom]{Exercices}
\item \hyperref[guide-section-phantom]{Guide bibliographique}
\item \hyperref[desirables-section-phantom]{Desiderata}
\item \hyperref[coding-section-phantom]{Style de programmation}
\item \hyperref[obsolete-section-phantom]{Obsolète}
\item \hyperref[fdl-section-phantom]{Licence de documentation libre GNU}
\item \hyperref[index-section-phantom]{Index généré automatiquement}
\end{enumerate}
\end{multicols}


\bibliography{my}
\bibliographystyle{amsalpha}

\end{document}
